\documentclass[a4paper,12pt]{article}

%\documentclass[a4paper,12pt, draft]{article}
\documentclass[a4paper,12pt]{article}

%%% Работа с русским языком
\usepackage{cmap}					% поиск в PDF
\usepackage{mathtext} 				% русские буквы в формулах
\usepackage[T2A]{fontenc}			% кодировка
\usepackage[utf8]{inputenc}			% кодировка исходного текста
\usepackage[english,russian]{babel}	% локализация и переносы
\usepackage{indentfirst}            % красная строка в первом абзаце
\frenchspacing                      % равные пробелы между словами и предложениями

%%% Дополнительная работа с математикой
\usepackage{amsmath,amsfonts,amssymb,amsthm,mathtools} % пакеты AMS
\usepackage{icomma}                                    % "Умная" запятая

%%% Свои символы и команды
\usepackage{centernot} % центрированное зачеркивание символа
\usepackage{stmaryrd}  % некоторые спецсимволы

\renewcommand{\epsilon}{\ensuremath{\varepsilon}}
\renewcommand{\phi}{\ensuremath{\varphi}}
\renewcommand{\kappa}{\ensuremath{\varkappa}}
\renewcommand{\le}{\ensuremath{\leqslant}}
\renewcommand{\leq}{\ensuremath{\leqslant}}
\renewcommand{\ge}{\ensuremath{\geqslant}}
\renewcommand{\geq}{\ensuremath{\geqslant}}
\renewcommand{\emptyset}{\ensuremath{\varnothing}}

\DeclareMathOperator{\sgn}{sgn}
\DeclareMathOperator{\ke}{Ker}
\DeclareMathOperator{\im}{Im}
\DeclareMathOperator{\re}{Re}

\newcommand{\N}{\mathbb{N}}
\newcommand{\Z}{\mathbb{Z}}
\newcommand{\Q}{\mathbb{Q}}
\newcommand{\R}{\mathbb{R}}
\newcommand{\Cm}{\mathbb{C}}
\newcommand{\F}{\mathbb{F}}
\newcommand{\id}{\mathrm{id}}
\newcommand{\imp}[2]{
	(#1\,\,$\ra$\,\,#2)\,\,
}
\newcommand{\System}[1]{
	\left\{\begin{aligned}#1\end{aligned}\right.
}
\newcommand{\Root}[2]{
	\left\{\!\sqrt[#1]{#2}\right\}
}
\newcommand{\RR}{\R}
\newcommand{\NN}{\N}
\renewcommand{\subseteq}{\subset}
\newcommand{\sub}{\subset}
\newcommand{\sconstr}{\;\vert\;}
\newcommand{\thus}{\implies}

\newcommand{\defeq}{\vcentcolon= }
\newcommand{\defev}{\stackrel{\Delta}{\Longleftrightarrow}}
\newcommand{\deriv}[3][1]{%
	\ifthenelse{#1>1}{%
		\frac{\dlta^{#1} {#2}}{\dlta {#3}^{#1}}
	}{%
		\frac{\dlta {#2}}{\dlta {#3}}
	}%
}

\renewcommand\labelitemi{$\triangleright$}

\let\bs\backslash
\let\lra\Leftrightarrow
\let\ra\Rightarrow
\let\la\Leftarrow
\let\emb\hookrightarrow

%%% Перенос знаков в формулах (по Львовскому)
\newcommand{\hm}[1]{#1\nobreak\discretionary{}{\hbox{$\mathsurround=0pt #1$}}{}}

%%% Работа с картинками
\usepackage{graphicx}    % Для вставки рисунков
\setlength\fboxsep{3pt}  % Отступ рамки \fbox{} от рисунка
\setlength\fboxrule{1pt} % Толщина линий рамки \fbox{}
\usepackage{wrapfig}     % Обтекание рисунков текстом

%%% Работа с таблицами
\usepackage{array,tabularx,tabulary,booktabs} % Дополнительная работа с таблицами
\usepackage{longtable}                        % Длинные таблицы
\usepackage{multirow}                         % Слияние строк в таблице

%%% Теоремы
\theoremstyle{plain}
\newtheorem{theorem}{Теорема}[section]
\newtheorem{lemma}{Лемма}[section]
\newtheorem{proposition}{Утверждение}[section]
\newtheorem*{exercise}{Упражнение}
\newtheorem*{problem}{Задача}

\theoremstyle{definition}
\newtheorem{definition}{Определение}[section]
\newtheorem*{corollary}{Следствие}
\newtheorem*{note}{Замечание}
\newtheorem*{reminder}{Напоминание}
\newtheorem*{example}{Пример}

\theoremstyle{remark}
\newtheorem*{solution}{Решение}

%%% Оформление страницы
\usepackage{extsizes}     % Возможность сделать 14-й шрифт
\usepackage{geometry}     % Простой способ задавать поля
\usepackage{setspace}     % Интерлиньяж
\usepackage{enumitem}     % Настройка окружений itemize и enumerate
\setlist{leftmargin=25pt} % Отступы в itemize и enumerate

\geometry{top=25mm}    % Поля сверху страницы
\geometry{bottom=30mm} % Поля снизу страницы
\geometry{left=20mm}   % Поля слева страницы
\geometry{right=20mm}  % Поля справа страницы

\setlength\parindent{15pt}        % Устанавливает длину красной строки 15pt
\linespread{1.3}                  % Коэффициент межстрочного интервала
%\setlength{\parskip}{0.5em}      % Вертикальный интервал между абзацами
%\setcounter{secnumdepth}{0}      % Отключение нумерации разделов
%\setcounter{section}{-1}         % Нумерация секций с нуля
\usepackage{multicol}			  % Для текста в нескольких колонках
\usepackage{soulutf8}             % Модификаторы начертания
\mathtoolsset{showonlyrefs=true} % показывать номера формул только у тех, у которых есть ссылки по eqref
%%% Содержаниие
\usepackage{tocloft}
\tocloftpagestyle{main}
%\setlength{\cftsecnumwidth}{2.3em}
%\renewcommand{\cftsecdotsep}{1}
%\renewcommand{\cftsecpresnum}{\hfill}
%\renewcommand{\cftsecaftersnum}{\quad}

%%% Шаблонная информация для титульного листа
\newcommand{\CourseName}{Математический анализ}
\newcommand{\FullCourseNameFirstPart}{\so{Введение в математический анализ}}
\newcommand{\SemesterNumber}{I}
\newcommand{\LecturerInitials}{Гусев Николай Анатольевич}
\newcommand{\CourseDate}{осень 2022}
\newcommand{\AuthorInitials}{Копанов Антон}
\newcommand{\AuthorInitialssecond}{Клищ Данил}
\newcommand{\VKLink}{https://vk.com/notnako}
\newcommand{\GithubLink}{https://github.com/MIPT-Group/Lectures_Tex_Club}
\newcommand{\VKLinksecond}{https://vk.com/dan.klishch}

%%% Колонтитулы
\usepackage{titleps}
\newpagestyle{main}{
	\setheadrule{0.4pt}
	\sethead{\CourseName}{}{\hyperlink{intro}{\;Назад к содержанию}}
	\setfootrule{0.4pt}                       
	\setfoot{ФПМИ МФТИ, \CourseDate}{}{\thepage} 
}
\pagestyle{main}  

%%% Нумерация уравнений
\makeatletter
\def\eqref{\@ifstar\@eqref\@@eqref}
\def\@eqref#1{\textup{\tagform@{\ref*{#1}}}}
\def\@@eqref#1{\textup{\tagform@{\ref{#1}}}}
\makeatother                      % \eqref* без гиперссылки
\numberwithin{equation}{section}  % Нумерация вида (номер_секции).(номер_уравнения)
\mathtoolsset{showonlyrefs= true} % Номера только у формул с \eqref{} в тексте.

%%% Гиперссылки
\usepackage{hyperref}
\usepackage[usenames,dvipsnames,svgnames,table,rgb]{xcolor}
\hypersetup{
	unicode=true,            % русские буквы в раздела PDF
	colorlinks=true,       	 % Цветные ссылки вместо ссылок в рамках
	linkcolor=black!15!blue, % Внутренние ссылки
	citecolor=green,         % Ссылки на библиографию
	filecolor=magenta,       % Ссылки на файлы
	urlcolor=NavyBlue,       % Ссылки на URL
}

%%% Графика
\usepackage{tikz}        % Графический пакет tikz
\usepackage{tikz-cd}     % Коммутативные диаграммы
\usepackage{tkz-euclide} % Геометрия
\usepackage{stackengine} % Многострочные тексты в картинках
\usetikzlibrary{angles, babel, quotes}

%%% Всю шаблонную информацию можно менять тут
\newcommand{\FullCourseNameFirstPart}{\so{КРАТНЫЕ ИНТЕГРАЛЫ И}}
\newcommand{\FullCourseNameSecondPart}{\so{ТЕОРИЯ ПОЛЯ}}
\newcommand{\SchoolName}{ФПМИ}
\newcommand{\TrackName}{ПМИ, ИВТ, КТ, РАНХиГС}
\newcommand{\SemesterNumber}{III}
\newcommand{\LecturerInitials}{Лукашов Алексей Леонидович}
\newcommand{\AutherInitialsFirst}{Климов Ярослав}
\newcommand{\AutherInitialsSecond}{Ильдаров Адам}
\newcommand{\AutherInitialsThird}{Евдокимов Егор}
\newcommand{\VKLinkFirst}{https://vk.com/klimov_yaroslav}
\newcommand{\VKLinkSecond}{https://vk.com/adamildarov}
\newcommand{\VKLinkThird}{https://vk.com/ea_evdokimov}
\newcommand{\OverleafLink}{https://www.overleaf.com/read/thzfwhhyghfn}
\newcommand{\GithubLink}{https://github.com/MIPT-Group/Lectures_Tex_Club}

\newcommand{\CourseName}{Кратные интегралы и теория поля} %  Используется в преамбуле для создания названия предмета в верхнем контитуле   
\newcommand{\CourseDate}{2020}           %  Используется в преамбуле для создания даты в нижнем контитуле и в title_page

%\includeonly{lectures/lect05,lectures/lect06}  % Чтобы скомпилировать только часть лекций

\begin{document}
    %%% Всю шаблонную информацию можно менять тут
\newcommand{\FullCourseNameFirstPart}{\so{ТЕОРИЯ ГРУПП}}
\newcommand{\SemesterNumber}{III}
\newcommand{\LecturerInitials}{Богданов Илья Игоревич}
\newcommand{\AutherInitials}{Даниил Дрябин}
\newcommand{\VKLink}{https://vk.com/ax_equals_b}
\newcommand{\GithubLink}{https://github.com/daniild71r/lectures_tex_club}

\begin{titlepage}
	\clearpage\thispagestyle{empty}
	\centering
	
	%\textit{Федеральное государственное автономное учреждение \\
	%	высшего профессионального образования}
	%\vspace{0.5ex}
	
	%\textbf{Московский Физико-Технический Институт \\ КЛУБ ТЕХА ЛЕКЦИЙ}
	\textbf{Московский физико-технический институт \\ Физтех-школа прикладной математики и информатики}
	\vspace{20ex}
	\vspace{13ex}
	
	{\textbf{\FullCourseNameFirstPart}}
	
	\SemesterNumber\ СЕМЕСТР  
	\vspace{1ex}
	
	Лектор: \textit{\LecturerInitials}
	
	\begin{figure}[!ht]
		\centering
		\includegraphics[width=0.4\textwidth]{images/logo_ltc.png}
	\end{figure}

	\begin{flushright}
		\noindent
		Автор 2020: \href{\VKLink}{\textit{\AutherInitials}}
		\\
		Дополнил 2022: \textit{Даниил Максимов}
		\\
		\href{\GithubLink}{\textit{Проект на Github}}
	\end{flushright}
	
	\vfill
	\CourseDate
	\pagebreak
\end{titlepage}

    \newpage
    \hypertarget{intro}{}
    \tableofcontents
    \newpage
    \section{Unix системы}

\subsection{Компоненты операционной системы}

\begin{Def}
	\underline{Ядро} --- программа, которая запускается самой первой. Имеет привилегированное положение на процессоре и может взаимодействовать с <<железом>> напрямую.
\end{Def}

\begin{Def}
	\underline{Базовые библиотеки}: зачастую входят в состав операционной системы, потому что привязаны к определенному ядру (версии ядра), чтобы максимально использовать его функциональность. 
\end{Def}

\begin{Def}
\underline{Служебные сервисы} --- ведение логов, обработка сетевых соединения, и т.д.
\end{Def}

\begin{Def}
\underline{Минимальная пользовательская оболочка} --- чтобы пользователь мог взаимодействовать с системой.
\end{Def}

\subsection{Процессы}

\begin{Def}
\underline{Процесс} --- любой экземпляр программы, работающей в операционной системе.
\end{Def}

Примеры:
\begin{itemize}
	\item Приложения
	\item Фоновые сервисы (демоны)
	\item Команды в терминале
\end{itemize}

Главные команды для мониторинга процессов:
\begin{itemize}
	\item ps -A -- список процессов
	\item pstree -- иерархия процессов
	\item top
\end{itemize}

\subsection{UNIX программы}

Программы:
\begin{itemize}
	\item \textbf{Бинарные исполняемые файлы.}
	\item \textbf{Текстовый файл, начинающийся с \#!.} После \#! указывается путь к интерпретатору (например, \#!/bin/bash). Это могут быть скрипты на Python/Perl/Escript или Shell скрипт.
\end{itemize}

Название файла и его расширение не имеет значения. Любой файл, имеющий атрибут 
<<исполняемый>> является программой.

\subsection{UNIX пользователи}

Единственные привилегированный пользователь: root (UID = 0).

Непривилегерованные пользователи:
\begin{itemize}
	\item Реальные пользователи, которые могут войти в систему (UID $\ge$ 1000)
	\item Фейковые пользователи, привязанные к сервисам (UID $<$ 1000).
\end{itemize}

\textbf{Зачем сервисам нужны фейковые пользователи?}
Допустим, у вас есть веб-сервер и сервер баз данных. Теоретически какой-то пользователь
может подключиться к веб-серверу и найти там уязвимость. Для того, чтобы 
минимизировать возможный ущерб от уязвимости в одном из процессов, процессы запускаются
под разными пользователями и не имеют права общаться друг с другом.

Непривилегерованные пользователи могут временно получить 
дополнительные привилегии:

\begin{itemize}
	\item su - запускает терминал от имени root
	\item sudo - запускает любую программу от имени root
	\item Запустить программу с атрибутом SUID и владельцем root. SUID атрибут означает, что при запуске программа
	исполняется от того пользователя, кто является владельцем файла. Кстати, так и 
	работает команда sudo, владелец sudo - root, а также у sudo проставлен атрибут SUID.
	\item Запустить программу с <<capabilities>> флагами. Эти флаги позволяют тонко настроить, что можно делать программе, а что - нельзя.
\end{itemize}

\subsection{UNIX межпроцессорное взаимодействие}

Процессы работают изолированно друг от друга. Единственный способ взаимодействовать друг с другом --- использовать методы межпроцессорного взаимодействия ядра (Inter-Process Communication, IPC):

\begin{itemize}
	\item {Сигналы} для коротких сообщений
	\item {Каналы и сокеты} для последовательных данных
	\item {Разделяемая память}, в которой хранятся большие куски данных
\end{itemize}

\subsection{Процесс запуска}

\begin{enumerate}
	\item {Загрузчик} с диска или UEFI материнской платы выполняется в привилегированном режиме
	\item Загрузчик запускает {ядро}, тоже в привилегированном режиме.
	\item Ядро запускает первый непривелигерованный процесс: {init или systemd}.
\end{enumerate}

\subsection{Сервисы (``Демоны'')}

\begin{Def}
	\underline{Сервисы} --- это специальные системные службы, которые работают в фоновом режиме.
\end{Def}

\newpage

Минимальный сет во многих системах:

\begin{itemize}
	\item \textbf{dhclient} -- держит активным полученный IP адрес в сети
	\item \textbf{getty} -- переключение между графическим/консольным входом в систему
	\item \textbf{sshd} -- подключение по ssh
	\item \textbf{ntpd или chronyd} -- синхронизация времени
	\item \textbf{syslogd} -- системный лог
\end{itemize}

\subsection{\#!/bin/sh}

В терминале вы работаете в интерпретаторе \texttt{shell}. \texttt{shell} предоставляет возможность выполнять
команды и последовательности команд.

Бывают разные интерпретаторы \texttt{shell}:
\begin{itemize}
	\item самодостаточный \texttt{shell} program: \textbf{FreeBSD}
	\item основной в \texttt{Linux}: \textbf{bash}
	\item debian: \textbf{dash}
	\item alpine: \textbf{busybox}
	\item maxOS X: \textbf{zsh}
\end{itemize}

\subsection{Системный демон (systemd)}
\begin{Def}
Процесс, который управляет всеми остальными процессами. При этом сам является обычным процессом. \texttt{Systemd} пришел на замену INIT-процессу.

\end{Def}

Чем systemd лучше?
\begin{itemize}
	\item Устранение интерпретации \texttt{shell} скриптов. 
	\item Запуск сервисов параллельно
\end{itemize}

\subsubsection{Концепты systemd}

\begin{itemize}
	\item Сервисы
	\item Цели:
		\begin{itemize}
			\item \underline{Стандартные цели}. Например, уровни \texttt{init} в System-V:
				\begin{itemize}
					\item poweroff.target (init 0)
					\item rescue.target (init 1)
					\item multi-user.target (init 3)
					\item graphical.target (init 5)
					\item reboot.target (init 6)
					\item sleep.target и hibertate.target (не присутствуют в System-V)
				\end{itemize}
			\item \underline{Специальные цели}
		\end{itemize}
\end{itemize}

\subsection{Фоновые задачи}

\begin{itemize}
	\item \textbf{Нет привязанного к задаче терминала.} Для того, чтобы сигнал \texttt{SIGHUP}
	не останавливал исполнение задачи используется команда \texttt{nohup}.
	\item \textbf{Пишут в лог.} Иногда пишут данные в какую-то старинную 
	локальную почтовую систему.
	\item \textbf{Запускается в единственном экземпляре.} Обычно это достигается
	с помощью file lock.
\end{itemize}

\subsection{Дистрибутивы \texttt{Linux}}

Ядро называется Linux. Базовая система --- минимальная юзабельная система.

\underline{Дистрибутив}:
\begin{itemize}
	\item Ядро
	\item Базовая система
	\item GUI оболочка
	\item Программное обеспечение 
\end{itemize}

\subsubsection{Какие есть дистрибутивы?} 

Общего назначения:
\begin{itemize}
	\item OpenSUSE
	\item Fedora
	\item Debian
	\item Ubuntu
\end{itemize}

Специфические дистрибутивы:
\begin{itemize}
	\item Alpine Linux --- очень легковесный, поэтому его
	хорошо ставить на виртуальную машину.
	\item Kali Linux --- используется для тестов на безопасность.
\end{itemize}


    \newpage
   	\begin{flushright}
    \textit{Лекция 2 (от 08.09)}
\end{flushright}
\section{$\S 3.$ Дифференцирование}
Пусть $z_0 \in \CC$, $B_r(z_0)$ и $f: B_r(z_0) \mapsto \CC$
\Def
Пусть $f: B_r(z_0) \mapsto \CC$. Если
\begin{equation}\label{(3.1)}
    \exists \dst \lim_{z \mapsto z_0} \frac{f(z)-f(z_0)}{z-z_0}
\end{equation}
то он называется \textbf{производной функции $f$ в точке $z_0$} (записывается
$f'(z_0)$).
\\
Можем расписать предел:
\begin{equation}\label{(3.2)}
    \forall \varepsilon > 0 \ \exists \delta(\varepsilon)>0: \ z \in B_\delta(z_0), \ \left| \frac{f(z) - f(z_0)}{z-z_0} - f'(z_0) \right| < \varepsilon
\end{equation}
\begin{equation}\label{(3.3)}
    f(z) = f(z_0) + f'(z_0)(z-z_0) + \alpha(z-z_0)
\end{equation}
\begin{equation}\label{(3.4)}
    \lim_{z \to z_0}\frac{\alpha(z-z_0)}{z-z_0} = 0
\end{equation}
\begin{align*}
  & f(z_0) = u(x,y) + iv(x,y)
\end{align*}
\Def
Говорят, что $f: B_r(z_0) \mapsto \CC$ \textbf{дифференцируема в точке $z_0$},
если
\begin{equation}\label{(3.5)}
    \exists A \in \CC: \ f(z) = f(z_0)+A(z-z_0)+o(z-z_0)
\end{equation}
\lemma
$f: B_r(z_0) \mapsto \CC$ дифференцируема в $z_0$ $\Leftrightarrow$ $\exists
f'(z_0)$, $A = f'(z_0)$
\pr
Очевидно.
\\
Обозначим: $\Delta x = x-x_0$, $\Delta y = y - y_0$, $\Delta z = z-z_0 = \Delta
x + i \Delta y$, $\Delta u = u(x,y)-u(x_0,y_0)$, $\Delta v = v(x,y)-v(x_0,y_0)$,
$\Delta f = f(x,y)-f(x_0,y_0) = \Delta u + i \Delta v$.
\theorem
$f: B_r(z_0) \mapsto \CC$ дифференцируема в $z_0$ т огда и только тогда, когда
\begin{itemize}
    \item $u(x,y)$, $v(x,y)$ дифференцируемы в $(x_0,y_0)$ (в $\RR^2$)
    \item выполняется условие Коши-Римана (УКР):
    \begin{equation}\label{(3.6)}
        \begin{cases}
            \dst \frac{\partial u}{\partial x} = \dst \frac{\partial v}{\partial y} \\
            \dst \frac{\partial v}{\partial x} = - \dst \frac{\partial u}{\partial y}
        \end{cases}
        % \left\{ \begin{matrix}
        %         \dst \frac{\partial u}{\partial x} = \dst \frac{\partial v}{\partial y} \\
        %         \dst \frac{\partial v}{\partial x} = - \dst \frac{\partial u}{\partial y}
        %     \end{matrix} \right.
    \end{equation}
\end{itemize}
При этом
\begin{equation}\label{(3.7)}
    f'(z_0) = \frac{\partial u}{\partial x}(x_0,y_0) + i\frac{\partial v}{\partial x}(x_0,y_0) = \frac{\partial v}{\partial y}(x_0,y_0) - i\frac{\partial u}{\partial y}(x_0,y_0)
\end{equation}
\pr
Докажем в обе стороны.
\begin{itemize}
    \item Необходимость.
    \begin{align*}
      & \exists f'(z_0) = a+ib = A \in \CC
    \end{align*}
    Значит, по определению дифференцируемости
    \begin{align*}
      & \Delta f = A \Delta z + \alpha(\Delta z); \ \alpha(\Delta z) = \alpha_1(\Delta x, \Delta y) + i\alpha_2(\Delta x, \Delta y)
    \end{align*}
    \begin{align*}
      & \left\{ \begin{matrix}
              \Delta u = a \Delta x - b \Delta y + \alpha_1(\Delta x, \Delta y) \\
              \Delta v = b \Delta x + a \Delta y + \alpha_2(\Delta x, \Delta y)
          \end{matrix} \right.
    \end{align*}
    \begin{align*}
      & \left| \alpha_1 \right| \leq \left| \alpha \right|, \ \left| \alpha_2 \right| \leq \left| \alpha \right|
    \end{align*}
    \begin{align*}
      & \lim_{(\Delta x; \Delta y) \to (0;0)}\left( \frac{\left| \alpha_1(\Delta x, \Delta y) \right|}{\sqrt{\left( \Delta x \right)^2 + \left( \Delta y \right)^2}} \right) \leq \lim_{\Delta z \to 0}\left( \frac{\left| \alpha(\Delta z) \right|}{\left| \Delta z \right|} \right) = 0
    \end{align*}
    Значит, $u$ дифференцируема, причем $\dst \frac{\partial u}{\partial x} =
    a$, $\dst \frac{\partial u}{\partial y} = -b$.
    \\
    Аналогично для $v$, причем $\dst \frac{\partial v}{\partial x} = b$, $\dst
    \frac{\partial v}{\partial y} = a$.
    \\
    Видим выполнимость УКР.
    \\
    Выполнимость \eqref{(3.7)} очевидна.
    \item Достаточность.
    \\
    Пусть $u$, $v$ дифференцируемы в $(x_0,y_0)$ и выполняется УКР. Тогда
    \begin{align*}
      & \Delta f = \Delta u + i \Delta v = \frac{\partial u}{\partial x} \Delta x + \frac{\partial u}{\partial y} \Delta y + \alpha_1(\Delta x, \Delta y) + i \left( \frac{\partial v}{\partial x} \Delta x + \frac{\partial v}{\partial y} \Delta y + \alpha_2(\Delta x, \Delta y)\right) = \\
      & \frac{\partial u}{\partial x} \Delta x - \frac{\partial v}{\partial x} \Delta y + \alpha_1(\Delta x, \Delta y) + i \left( \frac{\partial v}{\partial x} \Delta x + \frac{\partial u}{\partial x} \Delta y + \alpha_2(\Delta x, \Delta y)\right) = \left( \frac{\partial u}{\partial x} + i \frac{\partial v}{\partial x}\right)\cdot \\
      & \cdot \left( \Delta x + i \Delta y \right)+ \alpha_1(\Delta x, \Delta y) + i \alpha_2(\Delta x, \Delta y)
    \end{align*}
    Значит,
    \begin{align*}
      & \exists f'(z_0) = \frac{\partial u}{\partial x} + i \frac{\partial v}{\partial x}
    \end{align*}
\end{itemize}
\Example
\begin{align*}
  & w = z^2
\end{align*}
\begin{itemize}
    \item[I] Найдем производную по определению.
    \begin{align*}
      & \frac{\Delta w}{\Delta z} = \frac{\left( z_0-\Delta z \right)^2 - z_0^2}{\Delta z} = \frac{2z_0\Delta z + \left( \Delta z \right)^2}{\Delta z} = 2z_0 + \Delta z \us{\Delta z \to 0}{\to} 2z_0
    \end{align*}
    \item[II] Найдем производную с использованием теоремы.
    \begin{align*}
      & z^2 = x^2-y^2+2ixy \Rightarrow u = x^2-y^2, \ v = 2xy
    \end{align*}
    \begin{align*}
      & u_x = 2x = v_y; \ u_y = -2y = - v_x
    \end{align*}
    Видим выполнимость УКР, поэтому
    \begin{align*}
      & w' = u_x+iv_x = 2x + 2iy = 2z
    \end{align*}
\end{itemize}
\Example
\begin{align*}
  & w = \left| z \right|^2
\end{align*}
\begin{align*}
  & \left| z \right|^2 = x^2+y^2 \Rightarrow u = x^2+y^2, \ v = 0
\end{align*}
\begin{align*}
  & u_x = 2x \neq 0 = v_y; \ u_y = 2y \neq 0 = - v_x
\end{align*}
(за исключением точки $(0;0)$). Видим выполнение УКР только в точке $(0;0)$, а
значит, и функция дифференцируема только в этой точке (с производной, равной
$0$).
\Example
\begin{align*}
  & w = \ol{z}
\end{align*}
\begin{align*}
  & \ol{z} = x - iy \Rightarrow u = x, \ v = -y
\end{align*}
\begin{align*}
  & u_x = 1 \neq -1 = v_y; \ u_y = 0 = - v_x
\end{align*}
Видим, что УКР не выполняется ни в одной точке $\CC$, а значит, функция нигде не
дифференцируема.
\prop
Если $f(z)$, $g(z)$ дифференцируемы, то:
\begin{itemize}
    \item дифференцируема и сумма:
    \begin{align*}
      & \left( f+g \right)' = f'+g'
    \end{align*}
    \item дифференцируемо и произведение:
    \begin{align*}
      & \left( fg \right)' = f'g+fg'
    \end{align*}
    \item дифференцируемо и частное (когда $g(z) \neq 0$):
    \begin{align*}
      & \left( \frac{f}{g} \right)' = \frac{f'g+fg'}{g^2}
    \end{align*}
\end{itemize}
\pr
Доказательства легко провести через УКР с использованием действительного
анализа.
\section{$\S 4.$ Регулярные и гармонические функции}
\example
$B_r(z_0)$~--- односвязная область как в $\CC$, так и в $\CCC$.
\example
$K = \left\{ z: 0 < r_1 < \left| z \right| < r_2 < \infty \right\}$~--- не
односвязная область ни в $\CC$, ни в $\CCC$.
\example
$B_r(\infty)$~--- односвязная область в $\CCC$, но не в $\CC$.
\Def
Если $u:G \mapsto \RR^1$, $G \subseteq \RR^2$~--- область, $u \in C^2(G)$,
$\Delta u = 0$, где $\Delta = \dst \frac{\partial^2}{\partial x^2} + \dst
\frac{\partial^2}{\partial y^2}$, то $u$ называется \textbf{гармонической}.
\Def ~
\begin{itemize}
    \item[а)]$f: G \mapsto \CC$, где $G \subseteq \CC$~--- область, называется
    \textbf{регулярной} (аналитической, голоморфной), если $\forall z \in G \
    \exists f'(z)$ (если $f$ дифференцируема в каждой точке $G$).
    \item[б)]$f: G \mapsto \CC$, $G \subseteq \CC$ называется
    \textbf{регулярной в точке $z_0 \in G$}, если $\exists r > 0: B_r(z_0)
    \subseteq G$, $f$ регулярна на $B_r(z_0)$.
\end{itemize}
\Def
Функция, заданная на произвольном множестве, \textbf{регулярна} на нем, если она
регулярна в каждой его точке.
\Example
\begin{align*}
  & w = z^n, \ n \geq 2, \ n \in \NN
\end{align*}
Производная существует в любой точке:
\begin{align*}
  & \frac{\left( z+\Delta z \right)^n - z^n}{\Delta z} = \frac{nz^{n-1}\Delta z + \left( \Delta z \right)^2\left( P(z, \Delta z) \right)}{\Delta z} \us{\Delta z \to 0}{\to} nz^{n-1}
\end{align*}
Значит, функция всюду дифференцируема и регулярна на всей $\CC$.
\\
Аналогично можно доказать регулярность всех многочленов и рациональных функций.
\Def
\begin{align*}
  & e^z = e^x\left( \cos y + i \sin y \right), \ z = x+iy
\end{align*}
Функция периодична с периодом $2 i \pi$; ее компоненты $u = e^x \cos y$, $v =
e^x \sin y$.
Выполняются УКР:
\begin{align*}
  & \begin{cases}
      u_x = e^x \cos y = v_y \\
      v_x = e^x \sin y = -u_y
  \end{cases}
  & % \left\{ \begin{matrix}
    %       u_x = e^x \cos y = v_y \\
    %       v_x = e^x \sin y = -u_y
    %   \end{matrix} \right.
\end{align*}
Значит, $f'(z) = u_x + iv_x = u+iv = e^z$ (аналогично $(e^{iz})' = ie^{iz}$,
$(e^{-iz})' = -ie^{iz}$).
\\
Эти функции всюду дифференцируемы, а значит, и регулярны.
\Def
Тригонометрические функции:
\begin{align*}
  & \cos z = \frac{e^{iz} + e^{-iz}}{2}, \ \sin z = \frac{e^{iz} -e^{-iz}}{2i}
\end{align*}
Для них выполняются известные свойства тригонометрических функций:
\begin{align*}
  & \cos^2 z + \sin^2 z = 1
\end{align*}
\begin{align*}
  & \sin(z_1+z_2) = \sin z_1 \cos z_2 + \sin z_2 \cos z_1, \ \cos(z_1+z_2) = \cos z_1 \cos z_2 - \sin z_1 \sin z_2
\end{align*}
\begin{align*}
  & \left( \sin z \right)' = \cos z, \ \left( \cos z \right)' = - \sin z
\end{align*}
Но ограниченность не выполняется.
\\
Раскроем функцию через $x$ и $y$.
\begin{align*}
  & \sin z = \sin(x+iy) = \sin x \cos iy + \cos x \sin iy = \sin x \ch y + i \cos x \sh y \\
  & \Downarrow \\
  & u(x,y) = \sin x \ch y, \ v = \cos x \sh y
\end{align*}
Тогда
\begin{align*}
  & \left| \sin z \right| = \sqrt{\sin^2 x \ch^2 y + \cos^2 x \sh^2 y} = \sqrt{(1-\cos^2 x)\ch^2 y + (\ch^2 y - 1)\cos^2 x} = \\
  & = \sqrt{\ch^2 y - \cos^2 x} \us{y \to \infty}{\sim} \ch x \us{y \to \infty}{\to} \infty
\end{align*}
\theorem
Пусть $f: G \mapsto \CC$ регулярна в $G$, $f = u + iv$, $u,v \in C^2(G)$. Тогда
$u$ и $v$ являются гармоническими.
\pr
\begin{align*}
  & \Delta u = \frac{\partial}{\partial x}\left( \frac{\partial u}{\partial x} \right)+ \frac{\partial}{\partial y} \left( \frac{\partial u}{\partial y} \right) = \frac{\partial^2 v}{\partial x \partial y} - \frac{\partial^2 v}{\partial y \partial x} = 0
\end{align*}
Переход объясняется условием Коши-Римана для регулярной $f$.
\\
Аналогично можно доказать гармоничность $v$.
\\
Пара функций, связанных УКР, называется \textbf{сопряженными гармоническими}.
\theorem
Пусть $G$~--- односвязная область, $u: G \mapsto \RR$ гармоническая. Тогда
сущетвует регулярная $f$, такая, что $\Real f(z) = u(x,y)$.
\pr
Пусть
\begin{align*}
  & P(x,y) = -\frac{\partial u}{\partial y}, \ Q(x,y) = \frac{\partial u}{\partial x}
\end{align*}
Хотим найти $v$, такую, что
\begin{align*}
  & \begin{cases}
      \dst \frac{\partial v}{\partial x} = P(x,y) \\
      \dst \frac{\partial v}{\partial y} = Q(x,y) 
      \end{cases}
\end{align*}
(из УКР на $f$).
\begin{align*}
  & dv = Pdx+Qdy
\end{align*}
Для односвязной $G$ и потенциального векторного поля $(P,Q)$, т.~е. такого, что
выполняется условие:
\begin{align*}
  & \frac{\partial P}{\partial x} -\frac{\partial Q}{\partial x}' = - \frac{\partial^2 u}{\partial y^2} - \frac{\partial^2 u}{\partial x^2} = 0
\end{align*}
В таком случае
\begin{align*}
  & v(x,y) = \int_{(x_0,y_0)}^{(x,y)} P dx + Q dy + const
\end{align*}
Полученая функция $v$ гармоническая. тогда $f(z) = u+iv$ будет регулярной в
каждой точке $G$.
\\
Аналогично можно построить регулярную функцию $f$ по данной гармонической $v$.
\Example
Пусть
\begin{align*}
  & u = xy, \ \Delta u = 0
\end{align*}
Значит, $u$ гармоническая. УКР:
\begin{align*}
  & v_x = -u_y = -x \Rightarrow v(x,y) = \frac{-x^2}{2} + \varphi(y)
\end{align*}
\begin{align*}
  & v_y = u_x = y \Rightarrow v(x,y) = \frac{-x^2}{2} + \frac{y^2}{2} + C
\end{align*}
Значит,
\begin{align*}
  & \varphi'(y) = y = v_y
\end{align*}
Тогда
\begin{align*}
  & v(x,y) = \frac{y^2-x^2}{2} + C
\end{align*}
Значит,
\begin{align*}
  & f(x,y)= xy+i\left( \frac{y^2-x^2}{2} + C \right) = -\frac{i}{2} \left( x^2-y^2+2ixy \right) + iC = -\frac{iz}{2} + iC
\end{align*}

   	\newpage
   	\section{Целочисленная арифметика}

\subsection{Уровни абстракции <<железа>>}
\begin{itemize}
    \item <<p>> и <<n>> полупроводники, проводники, конденсаторы и резисторы
    \item Транзисторы и диоды
    \item Логические элементы
    \item Арифметическо-логические юниты
    \item Центральный процессор
    \item Компьютерная система
\end{itemize}

Компьютер состоит из элементарных блоков, из которых строятся электрические
элементы. Они в свою очередь образуют логические элементы.

Ключевая часть компьютера - центральный процессор. А главная часть ЦП - арифметическо-логический блок, который выполняет все вычисления.

\subsection{Типы сигналов}

\begin{Def}
    \underline{Сигналы} --- это изменения напряжения в электрической сети. Их можно регистрировать
и дальше обрабатывать.
\end{Def}

Сигналы бывают двух видов:
\begin{itemize}
    \item \textbf{Аналоговый.} Сила сигнала изменяется непрерывно с изменением напряжения.
    \item \textbf{Цифровой.} Есть всего два состояния - высокий, либо низкий.
\end{itemize}


\subsection{Дискретные сигналы}

\begin{itemize}
    \item \textbf{Логический <<0>>.} Все, что ниже определенного значения
        напряжения --- это нулевой уровень.
    \item \textbf{Логическая <<1>>.}
    \item \textbf{Неопределенное состояние.} Напряжение между 0 и 1, не приводящее к изменению состояния.
    \item \textbf{Hi-Z состояние.} Физически какой-то контакт внутри процессора отключается от общей схемы, чтобы не оказывать влияние на все остальные компоненты.
\end{itemize}

\subsection{Арифметико-логический юнит}

Базовый юнит центрального процессора для реализации операций над числами:
\begin{itemize}
	\item Битовая логика
	\item Побитовые сдвиги
	\item Сложение и вычитание
	\item Умножение и деление
\end{itemize}

 Управление ходом программы --- тоже арифметические операции. 
 
 \subsection{Представление чисел и переполнения}
 Было на семинарах...
 
 \subsection{Значения из нескольких байтов}
 
 \begin{itemize}
 	\item Минимальный адресуемый обьем данных --- 8 бит.
 	\item За один такт процессора он может адресовать одно \textbf{машинное слово} процессора - 32 или 64 бита.
 \end{itemize}

\subsection{Big-Endian и Little-Endian}

\begin{itemize}
	\item \textbf{Big-Endian:}
	Выглядит как человеческое представление байтов.
	сначала старшие байты, потом - младшие
	Пример: большинство RISC процессоров по умолчанию.
	\item \textbf{Little-Endian:}
	сначала младшие байты, потом старшие
	Пример: Процессоры Intel и ARM.
\end{itemize}

\subsection{Ковертирование endian-ов}

Беззнаковые типы можно конвертировать просто побитовыми сдвигами. Достаточно переставить байты в обратном порядке.

Со знаковыми типами есть проблемы, потому что там первый бит отвечает за знак.

\subsection{Упаковка структур}

Размер любой структуры должен быть кратек размеру машинного слова. Это связано с тем оптимизацией времени чтения/записи. 

Поэтому если в структурке хранится два поля - на 16 бит и на 8 бит, то стуктура будет занимать 32 бита, а не 16 + 8 = 24. Ведь размер машинного слова в x86 равен 32 битам.

Если же прописать у структуры специальный атрибут, компилятор уложит поля вплотную.

\subsection{Где используется атрибут packed?}

\begin{itemize}
	\item Бинарные данные
	\item Сетевые пакеты
	\item Драйвера
\end{itemize}

В остальных случаях использование packed ненужно и даже вредно.

   	\section{Команды процессора}

\subsection{Процессор}

Процессор умеет только манипулировать с целыми числами. Значение адреса - тоже число.

\subsection{Процессор с точки зрения пользователя}

У процессора есть:
\begin{itemize}
	\item Регистры - eax, ebx, ...
	\item Флаги - ZF, CF, SF
	\item Указатель на текущую команду (PC - pointer count)
\end{itemize}

\subsection{Команды x86}

nop - 0x00
halt - 0x10
call Dest - 0x80, Dest - 5 байтов
...

Команды кодируются переменным количеством байт. Количество байт может доходить до 6.

\subsection{Представление структур программы в виде простых инструкций}

Было на семинарах, на ассембелере писать уже должны уметь...

\subsection{CISC}

\begin{Def}
	\underline{CISC} --- \textbf{C}omplex \textbf{I}nstruction \textbf{S}et \textbf{C}omputing
\end{Def}

Процессоры: x86, Z80, PDP-11/VAX

Это тип архитектуры процессора, где:
\begin{itemize}
	\item Команд много, на все случаи жизни
	\item Кодируются переменным количеством байт
	\item Разные режимы адресации
	\item Упрощают жизнь программисту на ассемблере
\end{itemize}

Например, инструкция \textbf{loop} Address позволяет реализовать цикл одной командой.

\subsection{Конвейер}

\begin{itemize}
	\item Загрузка инструкции
	\item Декодирование команды
	\item Выполнение
	\item Доступ к памяти
	\item Записать результат в регистр
\end{itemize}

\subsection{Проблемы ковейеризации}

\begin{itemize}
	\item Инструкции имеют разную длину в байтах
	\item Разные инструкции выполняются за разное число тактов
	\item Работа с памятью и с регистрами задействует разные блоки процессора
\end{itemize}

Одна из задач компилятора с точки зрения оптимизации - минимизировать время простоя при конвейеризации. Для этого компилятор может менять порядок команд, если это не изменит результата, но поможет с ковейеризацией.

\subsection{RISC}

\begin{Def}
	\underline{RISC} --- \textbf{R}educed \textbf{I}nstruction \textbf{S}et \textbf{C}omputing
\end{Def}

Процессоры: ARM, AVR, MIPS, PowerPC

\begin{itemize}
	\item Только простейшие инструкции
	\item Инструкции фиксированной длины
	\item Почти у всех инструкций одинаковое время выполнения
	\item Адресация только по регистрам. Не можем читать напрямую из памяти
\end{itemize}

\subsection{Современные RICS процессоры}

\begin{itemize}
	\item \textbf{PowerPC:} Playstation, XBox, суперкомпьютеры IBM
	\item \textbf{MIPS:} WiFi/Bluetooth-адаптеры, процессоры в телевизорах
	\item \textbf{ARM:} Смартфоны
	\item \textbf{AVR:} Arduino
\end{itemize}

\subsection{Вещественная арифметика, IEEE 754}

Было на семинарах. Очень очень лень это все дублировать...




   	\section{Языки ассемблера и двоичный код}

\subsection{Стадии трансляции}

\begin{itemize}
	\item \textbf{Препроцессинг текста.} На выходе - текст.
	\item \textbf{Абстрактное дерево синтаксического разбора} и таблицы символов. Формат этого дерева внутренний, зависит от договоренностей.
	\item \textbf{Код на языке ассемблера или двоичный код.}
\end{itemize}

\subsection{Языки ассемблера}

\begin{itemize}
	\item Линейная последовательность команд
	\item Нет вложенных конструкций
	\item Нет вычислимых выражений. Точнее бывает, но только с константами и
	это обрабатывается на стадии препроцессинга (syntax sugar).
	\item Текст распознается регулярным выражением, а не контекстно-свободной грамматикой.
\end{itemize}

Любую программу языка ассемблера можно скомпилировать, а потом выполнить обратную 
операции - дизасемблирование. При этом код программы будет почти в точности совпадать. 

\subsection{Кодирование команд}

Есть две процессорные архитектуры:
\begin{itemize}
	\item \textbf{RISC.} Каждая команда длинной в машинное слово. Поэтому зная 
	начало кода можно извлечь любую инструкцию.
	\item \textbf{CISC.} Одна команда - произвольное количество байт.
\end{itemize}

\subsection{Команды и регистры}

\begin{itemize}
	\item Почти всегда процессор умеет выполнять действия только над регистрами. 
	\item Регистр - это не просто быстрая память. Это память, к которой процессор
	обращается напрямую.
	\item Доступ к памяти  - достаточно сложная операция.
	\item Некоторый процессоры (например, x86) имеют команды для адресации памяти, но на самом деле они сводятся к загрузке этой памяти в регистры.
\end{itemize}

\subsection{Кодирование команд}

В команде кодируется:
\begin{itemize}
	\item Какая именно команда
	\item C какими регистрами она работает
	\item Константы. При этом, не все константы можно закодировать
	в команде, потому что их слишком много. 
\end{itemize}

\subsection{Ключевая проблема}

Нельзя впихнуть невпихимуемое. Длина команды - всего 32 бита.

Как закодировать адрес? Очевидно, нельзя адресовать все виртуальное адресное пространство,
ведь это 4ГБ, а у нас есть всего 12 бит на кодирование адреса. \

Поэтому код команды устроен так: 8 бит на значение и 4 бита на сдвиг с вращением. С помощью этого можно закодировать наиболее часто используемые адреса.

\subsection{Procedure Linkage Table}

Библиотечные функции находятся далько от нашего коде в адресном пространстве. Поэтому
возникает проблема с тем, чтобы вызвать эту функцию, ведь как было показано выше, мы не
можем закодировать в команде произвольный адрес перехода. Эту проблему решает Procedure Linkage Table. 

\begin{Def}
	\underline{Procedure Linkage Table} --- это таблица в памяти, рядом с исполняемым кодом, где записаны реальные адреса библиотечных функций. Это позволяет вызвать библиотечную функцию, которая находится далеко в адресном пространстве.
\end{Def}

\subsection{Про кодирование команд}

У ARM есть расширения:
\begin{itemize}
	\item \textbf{Thumb.} 16 битные инструкции
	\item \textbf{Jazelle.} Декодирование байткода Java. Раньше было полезно, 
	чтобы запускать игрушки на слабых телефонах.
\end{itemize}

\subsection{Уровни абстракции}

\begin{itemize}
	\item Высокоуровневый язык
	\item Промежуточный язык
		\begin{itemize}
			\item байткод (Java, CLI или PyPy)
			\item LLVM - абстракция от ассемблеров разных архитектур
		\end{itemize}
	\item Язык ассемблера
	\item Двоичный код
\end{itemize}

\subsection{Зачем понимать кодирование команд}

\begin{itemize}
	\item Just-In-Time компиляция
	\item Эмулятор компьютерной системы
	\item Трансляция команд
\end{itemize}


	\section{Глава 11. Интегрирование дифференциальных форм. Теория поля.}
 \subsection{Элементы тензорной алгебры.}
Договоримся, что $E = \R^n, E^*$ --- сопряженное пространство, то есть пространство линейных функционалов (линейных форм, ковекторов). Вектора обозначаются $x_1, \ldots, x_p \in E$, $x_i=\sum\limits_{k=1}^n \xi_i^k e_k;$ формы обозначаются $y^1, \ldots, y^q\in E^*, y^j=\sum\limits_{l=1}^n \eta_l^j e^l$. 
\begin{Def}
	Полилинейной формой называется функция $U:E^p\times(E^*)^q\rightarrow \R$, линейная по каждому из своих аргументов, например, $U(\alpha x_1'+\beta x_1'', x_2, \ldots, x_p; y_1, \ldots, y^q)=\alpha U(x_1', x_2, \ldots, x_p;y^1, \ldots, y^q)+\beta U(x_1'', x_2, \ldots, x_p;y^1, \ldots, y^q), \alpha, \beta \in \R$.
\end{Def}
\begin{prop}
	Полилинейная форма однозначно определяется своими значениями на базисных элементах $E$ и $E^*$, то есть числами $w_{i_1, \ldots, i_p}^{j_1, \ldots, j_q} = U(e_{i_1}, \ldots, e_{i_p}; e^{j_1}, \ldots, e^{j_q})$, где $\{e_i\}_{i=1}^n$ --- базис $E$, $\{e^j\}_{j=1}^q$ --- двойственный базис $E^* (e^j(e_i)=\delta_i^j$ --- символ Кронекера).
\end{prop}
\begin{proof}
	\begin{multline*}
		U(x_1, \ldots, x_p; y^1, \ldots, y^q) = U(\sum\limits_{k_1=1}^n \xi_1^{k_1} e_{k_1}, \ldots, \sum\limits_{k_p=1}^n \xi_p^{k_p} e_{k_p}; \sum\limits_{l_1=1}^n \eta_{l_1}^1 e^{l_1}, \ldots, \sum\limits_{l_q=1}^n \eta_{l_q}^q e^{l_q})= \\ =
		\text{[пользуясь полилинейностью]} = \\ = \sum\limits_{k_1=1}^n \ldots \sum\limits_{k_p=1}^n\sum\limits_{l_1=1}^n\ldots\sum\limits_{l_q=1}^n \xi_1^{k_1}\ldots\xi_p^{k_p}\eta_{l_1}^1\ldots\eta_{l_q}^q U(e_{k_1}, \ldots, e_{k_p};e^{l_1}, \ldots, e^{l_q}) = \\ =
		\sum\limits_{k_1=1}^n \ldots \sum\limits_{k_p=1}^n\sum\limits_{l_1=1}^n\ldots\sum\limits_{l_q=1}^n \xi_1^{k_1}\ldots\xi_p^{k_p}\eta_{l_1}^1\ldots\eta_{l_q}^q w_{k_1, \ldots, k_p}^{l_1, \ldots, l_q} = \sum_{\overrightarrow{k}, \overrightarrow{l}}\xi^{\overrightarrow{k}}\eta_{\overrightarrow{l}}w_{\overrightarrow{k}}^{\overrightarrow{l}}.
	\end{multline*}
\end{proof}

\begin{Def}
	Набор чисел $\{w_{\overrightarrow{i}}^{\overrightarrow{j}}\}$ из утверждения 11.1.1 называется тензором полилинейной формы $U$.
\end{Def}

\begin{Def}
	Если $U, V$ --- формы, то $U+V, \alpha U$, где $\alpha\in\R$ определяются как $(U+~V)(x_1, \ldots, x_p; y^1, \ldots, y^q)=U(\ldots)+V(\ldots)$, и $(\alpha U)(\ldots)=\alpha(U(\ldots))$.
\end{Def}

\begin{prop}
	Множество полилинейных форм валентности $(p, q)$ образует линейное пространство $\Omega_p^q$.
\end{prop}

\begin{prop}
	Базис пространства $\Omega_p^q$ образуют формы $W_{j_1, \ldots, j_q}^{i_1, \ldots, i_p}$, задаваемые формулой $W_{j_1, \ldots, j_q}^{i_1, \ldots, i_p}(x_1, \ldots, x_p; y_1, \ldots, y^q)=\xi_1^{i_1}, \ldots, \xi_p^{i_p}\eta_{j_1}^1, \ldots, \eta_{j_q}^q$.
\end{prop}

\begin{Def}
	Тензорным произведением форм $U\in \Omega_{p_1}^{q_1}, V\in\Omega_{p_2}^{q_2}$ называется форма $U\otimes V\in \Omega_{p_1+p_2}^{q_1+q_2}$, задаваемая формулой
	\begin{multline*}
U\otimes V (x_1, \ldots, x_{p_1}, x_{p_1+1}, \ldots, x_{p_1+p_2}; y^1,\ldots, y^{q_1}, y^{q_1+1}, \ldots, y^{q_1+q_2})=\\=U(x_1, \ldots x_{p_1}; y^1, \ldots, y^{q_1})V(x_{p_1+1}, \ldots, x_{p_1+p_2}; y^{q_1+1}, \ldots, y^{q_1+q_2})
	\end{multline*} 
\end{Def}

\begin{prop}(Свойства тензорного призведения)
	\begin{enumerate}
		\item (Линейность) \\$(\alpha U_1+\beta U_2)\otimes V = \alpha(U_1\otimes V)+\beta(U_2\otimes V)$; \\
		$U\otimes (\alpha V_1+\beta V_2) = \alpha(U\otimes V_1)+\beta(U\otimes V_2), \alpha, \beta \in \R, U_1, U_2, U \in \Omega_{p_1}^{q_1}, V_1, V_2, V \in \Omega_{p_2}^{q_2}$
		\item (Ассоциативность) \\ $(U\otimes V)\otimes W = U\otimes (V\otimes W), U\in \Omega_{p_1}^{q_1}, V\in\Omega_{p_2}^{q_2}, W\in\Omega_{p_3}^{q_3}$
	\end{enumerate}
\end{prop}

\begin{prop}
	Базисные формы $W_{j_1, \ldots, j_q}^{i_1, \ldots, i_p} = e^{i_1}\otimes\ldots\otimes e^{i_p}\otimes e_{j_1}\otimes\ldots\otimes e_{j_q}$, где $e_j$ рассматриваются как элементы $(E^*)^*\cong E$
\end{prop}

Далее используем принцип Эйнштейна, согласно которому, будем опускать знак суммирования.

\textbf{Переход к другому базису.}

Пусть $\{e_i'\}_{i=1}^n$ --- новый базис $E$, \fbox{$e_i'=\alpha_i^je_j$} (здесь идет суммирование по $j$), $\{{f'}^j\}_{j=1}^n$ --- новый базис $E^*$, ${f'}^j= \beta_k^j f^k$, причем $\alpha_i^k \beta_k^j=\delta_i^j$. Получается, что $x=\xi^i e_i ={\xi'}^i e_i'$ --- вектор не зависит от того в каком базисе мы его записываем. Подставляя $e_i': x = \alpha_i^j{\xi'}^ie_j \Rightarrow \\ \xi^i = \alpha_i^j{\xi'}^i \Leftrightarrow {\xi'}^i=\beta_j^i\xi^j$. Для форм получаем: \fbox{$\eta_j'=\alpha_j^k\eta_k$}. Заметим, что базис в $E$ преобразуется так же как и координаты форм (выделены в рамках) --- это называется ковариантный закон. Если данное условие не выполняется, то выполняется контравариантный закон.
\begin{prop}
	Тензоры полилинейной формы принадлежащей $\Omega_p^q$ преобразуются при замене двойтсвенных базисов $p$ раз ковариантно и $q$ раз контравариантно, то есть ${w'}_{i_1, \ldots, i_p}^{j_1, \ldots, j_q}=\underbrace{\sum}_{p+q\text{ раз}}\beta^{j_1}_{l_1}\ldots\beta_{l_q}^{j_q}w_{k_1, \ldots, k_p}^{l_1, \ldots, l_q}\alpha_{i_1}^{k_1}\ldots\alpha_{i_p}^{k_p}$
\end{prop}

\begin{proof} По определению тензора формы:
	$${w'}_{i_1, \ldots, i_p}^{j_1, \ldots, j_q}=U({e'}_{i_1}, \ldots, {e'}_{i_p}; {f'}^{j_1}, \ldots, {f'}^{j_q}) = \alpha_{i_1}^{k_1}\ldots \alpha_{i_p}^{k_p}\beta_{l_1}^{j_1}\ldots\beta_{l_q}^{j_q}U(e_{k_1}, \ldots, e_{k_p};f^{l_1}, \ldots, f^{l_q})$$
\end{proof}

\begin{prop}
	Набор чисел $\{{w}_{i_1, \ldots, i_p}^{j_1, \ldots, j_q}\}$, $p$ раз ковариантный и $q$ раз контравариантный задает тензор полилинейной формы $U\in \Omega_p^q$ по формуле: $$U={w}_{i_1, \ldots, i_p}^{j_1, \ldots, j_q}f^{i_1}\otimes\ldots\otimes f^{i_p}\otimes e_{j_1}\otimes\ldots\otimes e_{j_q}$$
\end{prop}

\begin{proof}
	Пусть задана полилинейная форма $U(x_1, \ldots, x_p; y_1, \ldots, y^q)$, хотим доказать, что если мы переходим к другому базису, то значение формы будет тем же самым, так как векторы и ковекторы не зависят от базиса. Но эту форму мы определяли неким законом, который зависит от базиса. А теперь говорим что эта форма от базиса не зависит. Значит это надо проверить. Временно обозначим $\widetilde{U} = {w'}_{i_1, \ldots, i_p}^{j_1, \ldots, j_q}{f'}^{i_1}\otimes\ldots\otimes {f'}^{i_p}\otimes {e'}_{j_1}\otimes\ldots\otimes {e'}_{j_q}$ в измененном базисе. Тогда по предыдущему утверждению: $$\widetilde{U} = \beta_{l_1}^{j_1}\ldots\beta_{l_q}^{j_q}{w}_{k_1, \ldots, k_p}^{l_1, \ldots, l_q}\alpha_{i_1}^{k_1}\ldots\alpha_{i_p}^{k_p}\cdot\beta^{i_1}_{r_1}f^{r_1}\otimes\ldots\otimes\beta^{i_p}_{r_p}f^{r_p}\otimes\alpha_{j_1}^{s_1}e_{s_1}\otimes\ldots\otimes\alpha_{j_q}^{s_q}e_{s_q}.$$
	
	В этом выражении присутствуют такие конструкции: $\sum\limits_{i_1}\alpha_{i_1}^{k_1}\beta_{r_1}^{i_1} = \delta_{r_1}^{k_1}$. Тогда\\ $\widetilde{U} = w_{k_1\ldots k_p}^{l_1\ldots l_q}\delta_{r_1}^{k_1}\ldots\delta_{r_p}^{k_p}\delta_{l_1}^{s_1}\ldots\delta_{l_q}^{s_q}f^{r_1}\otimes\ldots\otimes f^{r_p}\otimes e_{s_1}\otimes\ldots\otimes e_{s_q}$. Суммирование символов Кронекера по $r$ и по $s$ в итоге даст $\widetilde{U} = w_{k_1\ldots k_p}^{l_1\ldots l_q}f^{k_1}\otimes\ldots\otimes f^{k_p}\otimes e_{s_1}\otimes\ldots\otimes e_{s_q}$.
\end{proof}

\textbf{Примеры}
\begin{enumerate}
	\item Вектор $x\in E$ с координатами $(\xi^1, \ldots, \xi^n)$ --- тензор валентности $(0, 1)$.
	
	Ковектор $y\in E^*$ с координатами $(\eta^1, \ldots, \eta^n)$ --- тензор валентности $(1, 0)$.
	\item Матрица линейного оператора $T: y= Tx, \eta^j=\tau_i^j\xi^i$, тензор $\{\tau_i^j\}$ валентности $(1,1)$, так как ${\eta'}^i=\beta_i^j\eta^i=\beta_i^j\tau_k^i\xi^k=\underbrace{\beta_i^j\tau^i_k\alpha_l^k}_{{\tau'}_l^j}{\xi'}^l$
\end{enumerate}

\begin{Def}
	$W\in\Omega_p^0$ называется симметрической, если она не изменяется при любой пререстановке ее аргументов.
\end{Def}

\begin{Def}
	$W\in\Omega_p^0$ называется антисимметрической или кососиммметрической, если при любой транспозиции пары ее аргументов, она меняет знак.
\end{Def}

\begin{prop}
	$W$ --- антисимметрическая $\Leftrightarrow$ она равна нулю, если какие-то два ее аргумента совпадают.
\end{prop}

\begin{proof}
	Пусть $W$ --- антисимметрическая, $x_k=x_l=a$. Тогда $W(x_1, \ldots, x_k, \ldots, x_l, \ldots, x_p)=-W(x_1, \ldots, x_k, \ldots, x_l, \ldots, x_p) = 0$.
	
	Пусть $W$ --- такая форма, что если два ее аргумента совпадают,то она нуль. 
	\begin{multline*}
		0=
		W(x_1, \ldots, x_k+x_l, \ldots, x_k+x_l, \ldots, x_p)=\\=
		\underbrace{W(x_1, \ldots, x_k, \ldots, x_k, \ldots, x_p)}_{=0}+
		W(x_1, \ldots, x_k, \ldots, x_l, \ldots, x_p)+\\+
		W(x_1, \ldots, x_l, \ldots, x_k, \ldots, x_p)+
		\underbrace{W(x_1, \ldots, x_l, \ldots, x_l, \ldots, x_p)}_{=0}
	\end{multline*}
Оставшиеся два слагаемых отличаются положением двух аргументов и в сумме равны 0, что и треюовалось доказать.
\end{proof}

\begin{corollary}[1]
	Если $x_1, \ldots, x_p$ --- линейно зависимы, то антисимметрическая $W(x_1, \ldots, x_p)=0$.
\end{corollary}

\begin{corollary}[2]
	Если $p>n$, то антисимметрическая $W(x_1, \ldots, x_p)\equiv0$
\end{corollary}

\begin{Def}
	Пусть $\pi_p=(i_1, \ldots, i_p)$ --- перестановка индексов $\{1, 2, \ldots, p\}$. Действие перестановки на форму $\pi_pW$ --- это форма. $\pi_pW(x_1, \ldots, x_p):=W(x_{i_1}, \ldots, x_{i_p})$
\end{Def}

\begin{prop}
	Если $W$ --- симметрическая, то $\pi_pW=W$ для $\forall \pi_p\in S_p$, где $S_p$ --- группа перестановок. (читать: для любой перестановки).
\end{prop}

\begin{prop}
	Если $W$ --- антисимметрическая, то $\pi_pW=\sgn(\pi_p)W$
\end{prop}

\begin{prop}
	$\forall W\in \Omega_p^0, \dfrac{1}{p!}\sum\limits_{\pi_p\in S_p}\pi_pW$ --- симметрическая, а $\dfrac{1}{p!}\sum\limits_{\pi_p\in S_p}\sgn(\pi_p)\pi_pW$ --- антисимметрическая
\end{prop}

\begin{proof}
	Пусть $\pi_p'$ --- некоторая транспозиция, подействуем ею на нашу сумму: 
	\begin{multline*}\pi_p'\left(\dfrac{1}{p!}\sum\limits_{\pi_p\in S_p}\sgn(\pi_p)\pi_pW\right)=\dfrac{1}{p!}\sum\limits_{\pi_p\in S_p}\sgn(\pi_p)\pi'_p\pi_pW = \\ =
	\dfrac{1}{p!}\sum\limits_{\pi_p\in S_p}\sgn^2(\pi'_p)\sgn(\pi_p)\pi'_p\pi_pW = \sgn(\pi'_p)\dfrac{1}{p!}\sum\limits_{\pi_p\in S_p}\sgn(\pi'_p\circ\pi_p)(\pi'_p\circ\pi_p)W, 
\end{multline*}
	так как $\sgn(\pi'_p\circ\pi_p)=\sgn(\pi'_p)\cdot\sgn(\pi_p)$. Суммируя по всем перестановкам, получаем исходную форму, со знаком впереди.
\end{proof}

\begin{Def}
	Симметризация формы $W\in\Omega_p^0: \sym W =\dfrac{1}{p!}\sum\limits_{\pi_p\in S_p}\pi_pW$. Антисимметризация формы $W\in\Omega_p^0: \asym W =\dfrac{1}{p!}\sum\limits_{\pi_p\in S_p}\sgn(\pi_p)\pi_pW$
\end{Def}

\begin{prop}
	$\sym W$ --- симметрическая форма, $\asym W$ --- антисимметрическая форма. $\sym , \asym $ --- линейные операторы на $\Omega_p^0$, сохраняющие симметрические и антисимметрические формы соответственно.
\end{prop}

\begin{Def}
	Если $U\in\Omega_p^0, V\in\Omega_q^0$ --- антисимметрические, то их внешним произведением называется $U\wedge V = \dfrac{(p+q)!}{p!q!}\asym(U\otimes V)$
\end{Def}


\begin{lemma}
	Если $U\in \Omega_p^0, V\in\Omega_q^0$, то $\asym(\asym U\otimes V)=\asym(U \otimes \asym V)=asym(U\otimes V)$
\end{lemma}

\begin{proof}Распишем:\\
	\mbox{$T = \asym(\asym U\otimes V)=\asym\left(\left(\dfrac{1}{p!}\sum\limits_{\pi\in S_p}\sgn(\pi)\pi U\right)\otimes V\right) = \dfrac{1}{p!}\sum\limits_{\pi\in S_p}\sgn(\pi)\asym((\pi U)\otimes V)$}. Рассмотрим отдельно на произвольном наборе векторов: 
	\begin{multline*}
		(\pi U)\otimes V(x_1, \ldots, x_{p+q}) = (\pi U)(x_1, \ldots, x_p)V(x_{p+1}, \ldots, x_{p+q}) =\\= U(x_{\pi(1)}, \ldots, x_{\pi(p)})V(x_{p+1}, \ldots, x_{p+q}) = \varpi_\pi (U\otimes V)(x_1, \ldots, x_{p+q}),
	\end{multline*}
	где $\varpi_\pi(i)=\pi(i), i=1, \ldots, p; \varpi_\pi(i)=i, i=p+1, \ldots, p+q$. 
	
	Тогда $T = \dfrac{1}{p!}\sum\limits_{\pi\in S_p}\sgn(\pi)\asym(\varpi_\pi(U\otimes V))\overset{\text{утв. 10}}{=} \dfrac{1}{p!}\sum\limits_{\pi\in S_p}\sgn(\pi)\sgn(\varpi_\pi)\asym(U\otimes V)$, так как  $\sgn(\pi)=\sgn(\varpi_\pi)$, то $T = \asym(U\otimes V)$.
\end{proof}

Введем обозначение $\Lambda_p$ --- линейное пространство антисимметрических форм из $\Omega_p^0$

\begin{theorem}(Основные свойства внешнего произведения)
	\begin{enumerate}
		\item (Дистрибутивность)\\$(\alpha U_1+\beta U_2)\wedge V = \alpha(U_1\wedge V)+\beta(U_2\wedge V); U\wedge(\alpha V_1+\beta V_2)=\alpha( U\wedge V_1)+\beta (U\wedge V_2),\\ (\forall \alpha, \beta \in \R), U_1, U_2, U\in \Lambda_p, V_1, V_2, V\in \Lambda_q$
		\item (Ассоциативность)\\$(U\wedge V)\wedge W = U \wedge (V\wedge W) = \dfrac{(p+q+r)!}{p!q!r!}\asym(U\otimes V\otimes W),\\ 
		U\in\Lambda_p, V\in\Lambda_q, W\in\Lambda_r$
		\item (Антикоммутативность) \\ $U\wedge V = (-1)^{pq}V\wedge U, U\in\Lambda_p, V\in\Lambda_q$
	\end{enumerate}
\end{theorem}

\begin{proof}
	\begin{enumerate}
		\item Следует из свойства 1 утверждения 4, определения, а также утверждения 12.
		
		\item \begin{multline*}
			(U\wedge V)\wedge W = \dfrac{((p+q)+r)!}{(p+q)!r!}\asym((U\wedge V)\otimes W)=\\ =\dfrac{(p+q+r)!}{(p+q)!r!}\asym\left(\dfrac{(p+q)!}{p!q!}\asym(U\otimes V)\otimes W\right) =\\= \dfrac{(p+q+r)!}{p!q!r!}\asym(\asym(U\otimes V)\otimes W)=\\=\text{[по лемме 1]}=\dfrac{(p+q+r)!}{p!q!r!}\asym((U\otimes V)\otimes W).
		\end{multline*}
	
		\item 
		$U\wedge V = \dfrac{(p+q)!}{p!q!}\asym(U\otimes V)=\dfrac{1}{p!q!}\sum\limits_{\pi\in S_{p+q}}\sgn(\pi)(\pi (U\otimes V))$.
		
		 Отдельно рассмотрим 
		 $T = \pi(U\otimes V)(x_1, \ldots, x_{p+q})=U\otimes V(x_{\pi(1)}, \ldots, x_{\pi(p+q)})=\\=U(x_{\pi(1)}, \ldots, x_{\pi(p)})V(x_{\pi(p+1)}, \ldots, x_{\pi(p+q)})$. Введем новую перестановку:
		 
		  $\pi'(i)=\pi(p+1), i=1, \ldots, q, \pi'(i)=\pi(i-q), i=q+1, \ldots, p+q$. 
		Тогда $T = U(x_{\pi'(q+1)}, \ldots, x_{\pi'(q+p)})V(x_{\pi'(1)}, \ldots, x_{\pi'(q)})=\pi'V\otimes U(x_1, \ldots, x_{p+q})$. 
		
		Вернемся к выражению $U\wedge V = \dfrac{1}{p!q!}\sum\limits_{\pi\in S_{p+q}}\sgn(\pi)(\pi'(V\otimes U))=\dfrac{(-1)^{pq}}{p!q!}\sum\limits_{\pi'\in S_{p+q}}\sgn(\pi')(\pi'(V\otimes U)) = (-1)^{pq}\dfrac{(p+q)!}{p!q!}\asym(V\otimes U)=(-1)^{pq}V\wedge U$, так как знаки перестановок $\pi$ и $\pi'$ отличаюстя на $(-1)^{pq}$.
	\end{enumerate}
\end{proof}
















	\section{Аллокация памяти}

\subsection{Физическая память}

Основные параметры:
\begin{itemize}
	\item частота шины (например, 1500 МГц).
	\item передача за одну секунду (например: 3000 MT/s).
\end{itemize}

\begin{Def}
	\underline{Латентность} --- задержка между командой и передачей данных.
\end{Def}

\subsection{Физическое адресное пространство}

Диапазон --- $2^{bits\_count}$

\subsection{Адресное пространство процесса}

У каждого процесса есть свое адресное пространство. Оно никак не пересекается
с адресным пространством других процессов.

Ограничения: 
\begin{itemize}
	\item 32-битные системы --- 4 Гб.
	\item 64-битные --- 256 Тб.
\end{itemize}

Пример содержимого:
\begin{itemize}
	\item Пространство ядра --- 0xc0000000
	\item Стек --- 0xbf800000
	\item Разделяемые библиотеки
	\item Куча
	\item Программа --- 0x004d9000
\end{itemize}

Для того, чтобы реализовать менеджмент памяти используется \textbf{страничная адресация}.

\subsection{Страничная адресация}

\begin{itemize}
	\item Разбиение памяти на \textbf{страницы} фиксированного размера.
	\item У каждой страницы есть свои атрибуты.
\end{itemize}

\subsection{Страничная адресация в x86-32}

32-битный адрес в виртуальном адресном пространстве разбивается на 3 части --- 
индекс в таблице первого уровня (10 бит), индекс в таблице второго уровня (10 бит) и 
смещение внутри страницы (12 бит).

\begin{itemize}
	\item Специальный регистр CR2 указывает на таблицу первого уровня текущего процесса.
	\item Каждая запись в таблицу первого уровня указывает на таблицу второго уровня.
	\item Записи в таблице второго уровня содержат указатели на страницы.
\end{itemize}

\begin{itemize}
	\item Размер страницы --- 4 Кб
	\item В каждой таблице есть до 1024 записей
\end{itemize}

Советую посмотреть \href{https://youtu.be/uSkrfUIYhqw?t=1355}{участок}
30 сек для лучшего понимания.

При кодировании у нас остается 12 бит для флагов --- это настройки страницы.

\subsection{Нарушение доступа к странице}

Флаг P (present --- файл существует в физицеской памяти) равен 0. Тогда процессор бросает исключение. Ядро обрабатывает это исключение.

\subsection{Аппаратная поддержка}

\begin{itemize}
	\item MMU (Memory Managment Unit) считает реальный указатель на используя указатель
	в виртуальном адресном пространстве.
	\item TLB (Translation Lookaside Buffer) предназначен для хранения текущих 
	таблиц в памяти.
	\item Размер страниц прописан на уровне процессора, вы не можете его поменять программно.
\end{itemize}

\subsection{Аллокаторы}

\begin{itemize}
	\item C++: стандартные new и new[] используют malloc/calloc. Это можно настраивать.
	\item C: malloc/calloc.
\end{itemize}

\subsection{Выделение памяти на кучу}

За счет чего реализованы библиотечные функции malloc и calloc?

Для этого используется один из системных вызовов:
\begin{itemize}
	\item \textbf{brk} --- устаревшая, но простая. Используется в glibc malloc до 128 Кб.
	\item \textbf{mmap} --- более гибкий инструмент.
\end{itemize}

\subsection{mmap}

Было на семах. 

Узнать больше: \textit{man mmap}

\subsection{Ограничения на память}

Ограничения можно посмотреть с помощью \textit{ulimit -a}.

\subsection{Что если память кончилась?}

Out-of-memory убийца процессов:
\begin{itemize}
	\item Найти жертву для убийства
	\item Очистить память
\end{itemize}

С каждый процессом связан некоторый рейтинг (/proc/[pid]/oom\_score), на основе которого
выбирается жертва.

Рейтинг вычисляется ядром. У процессов с высоким потреблением памяти рейтинг более
высокий, у привилегированных процесоов --- более маленький.

\subsection{Overcommit}

\begin{Def}
	\underline{Overcommit} --- это когда мы выделяем памяти больше, чем есть в системе.
	Это не вылетит пока мы не захотим ее использовать.
\end{Def}

Настроить стратегию при overcommit можно командой sysctl. 

Варианты стратегий:
\begin{itemize}
	\item OVERCOMMIT\_ALWAYS
	\item OVERCOMMIT\_GUESS --- стоит по умолчанию. Решение выносится на
	усмотрение ядра. Почти всегда это OVERCOMMIT\_ALWAYS.
	\item OVERCOMMIT\_NEVER
\end{itemize}
	\begin{flushright}
    \textit{Лекция 11 (от 12.10)}
\end{flushright}
\Def
Пусть в области $G$ задано семейство гладких кривых $\left\{ \gamma_\alpha
\right\}$, $\alpha \in [a,b]$, заданных при помощи параметра $t$: $z(t,
\alpha)$, $t \in [0;1]$. Если $z(t, \alpha)$ и $z'_t(t, \alpha)$ непрерывны на
$[0;1] \times [a,b]$, то говорят, что в области $G$ $\gamma_{\alpha}$
\textbf{задает непрерывную деформацию кривой $\gamma_a$ в кривую $\gamma_b$}.
\\
Если $\forall \alpha \in [a;b] \ z(0, \alpha) = z_0, \ z(1, \alpha) = z_1$, то
это называют \textbf{гомотетией $\gamma_a$ в $\gamma_b$}.
\theorem
Пусть в $\CC \setminus \{0\}$ задана непрерывная деформация $\gamma_a$ в
$\gamma_b$ через семейство $\left\{ \gamma_\alpha \right\}_{\alpha \in [a;b]}$.
Пусть $\exists A \in \CC, \ A \neq 0$ такое, что $\forall \alpha \in [a;b] \
z(1,\alpha) = Az(0, \alpha)$ и $I(\alpha) = \Delta_{\gamma_\alpha}\argt z =
\Delta_{[0;1]}\argt z(t, \alpha)$. Тогда $\forall \alpha \in [a,b] \ I(\alpha) =
const$. В частности, $\Delta_{\gamma_a}\argt z = \Delta_{\gamma_b}\argt z$.
\pr
По теореме $14.1$ $\forall \alpha \in [a;b] \ \exists \varphi(t, \alpha)$,
непрерывная по $t \in [0;1]$. $\varphi(t, \alpha) \in \Arg z (t, \alpha)$.
\begin{align*}
  & I(\alpha) = \Delta_{\gamma_\alpha}\argt z = \varphi(1, \alpha) - \varphi(0, \alpha) \in \Arg \frac{z(1, \alpha)}{z(0, \alpha)}
\end{align*}
\begin{align*}
  & \frac{z(1, \alpha)}{z(0, \alpha)} = A \neq 0 \Rightarrow \forall \psi_0 \in \Arg A \hookrightarrow I(\alpha) = \psi_0+2\pi k(\alpha), \ k(\alpha) \in \ZZ
\end{align*}
Значит, $I(\alpha)$ ступенчатая; но по определению приращения аргумента через
интеграл она непрерывна. Значит, она постоянна.
\\
$\forall z(t) \in C[0;1]$ таких, что $\forall t \in [0;1] \ z(t) \neq 0$ и $r =
\dst \min_{t \in [0;1]} \left| z(t) \right|$ положим $\varepsilon\in \left( 0;
    \dst \frac{r}{2} \right)$. Тогда существует \textbf{гладкая
  $\varepsilon$-аппроксимация $z(t)$}: $\forall t \in [0;1] \ z_\varepsilon(t)
\neq 0$, $\left| z(t) - z_\varepsilon(t) \right| \leq \varepsilon$,
$z_\varepsilon(0) = z(0)$, $z_\varepsilon(1) = z(1)$.
\\
Докажем это, построив пример. Действительно, по теореме Вейерштрасса
\begin{align*}
  & \exists P_n(t): \ \forall t \in [0;1] \left| z(t) - P_n(t) \right| < \frac{\varepsilon}{2}
\end{align*}
Положим
\begin{align*}
  & z_\varepsilon(t) = P_n(t) + \left( z(0)-P_n(0) \right)\left( 1-t \right) + \left( z(1)-P_n(1) \right)t
\end{align*}
Для такой функции выполняется $z(0) = z_\varepsilon(0)$, $z(1) =
z_\varepsilon(1)$; также
\begin{align*}
  & \abs{z(t) - z_\varepsilon(t)} \leq \abs{z(t) - P_n(t)} + \left| z(0)-P_n(0) \right| \left( 1-t \right) + \left| z(1)-P_n(1) \right|t < \frac{\varepsilon}{2} + \frac{\varepsilon}{2}(1-t) + \frac{\varepsilon}{2}(t) = \varepsilon
\end{align*}
и, наконец,
\begin{align*}
  & \abs{z_\varepsilon(t)} \geq \left| z(t) \right| - \left| z(t) - z_\varepsilon(t) \right| \geq r - \varepsilon \geq \frac{r}{2} > 0
\end{align*}
\Def
\textbf{Приращением аргумента непрерывной функции $z(t)$ вдоль непрерывной
  кривой $\gamma$ на $[0;1]$} называется $\Delta_{[0,1]}\argt z(t) =
\Delta_{[0,1]}\argt z_\varepsilon(t)$ для любой $\varepsilon$-аппроксимации
$z(t)$ при малых $\varepsilon$ ($\varepsilon \in \left( 0, \dst \frac{r}{2}
\right)$).
\prop
Определение $[14.4]$ не зависит от выбора $z_\varepsilon$.
\pr
Пусть $z_\varepsilon$~--- гладкая $\varepsilon$-аппроксимация,
$\tilde{z}_\varepsilon$~--- другая гладкая $\varepsilon$-аппорксимация.
Тогда
\begin{align*}
  & z(t, \alpha) = \alpha z_\varepsilon(t) + (1-\alpha)\tilde{z}_\varepsilon(t), \ \alpha \in [0;1], \ t \in [0;1]
\end{align*}
\begin{align*}
  & \abs{z(t, \alpha)} \geq \left| z(t) \right| - \left| z(t, \alpha) - z(t) \right| > r - \varepsilon \geq \frac{r}{2} > 0
\end{align*}
Значит, задана непрерывная деформация, и по теореме $14.3$ получаем, что
\begin{align*}
  & \Delta_{[0;1]}\argt z_\varepsilon(t) = \Delta_{[0;1]}\argt \tilde{z}_\varepsilon(t)
\end{align*}
\corollary
Пусть $\os{\circ}{\gamma}$~--- замкнутая непрерывная кривая, $0 \not \in
\os{\circ}{\gamma}$; тогда $\exists k(\os{\circ}{\gamma})\in \ZZ$:
$\Delta_{\os{\circ}{\gamma}}\argt z = 2 \pi k(\os{\circ}{\gamma})$.
\section{$\S 15.$ Регулярные ветви многозначных функций $\{\sqrt[n]{z}\}$ и $\Ln
  z$.}
Определим на $\CC\setminus\{0\}$:
\begin{equation}\label{(15.1)}
    \left\{ \sqrt[n]{z} \right\} = \sqrt[n]{\left| z \right|}\exp\left( \frac{i}{n} \Arg z \right)
\end{equation}
\begin{equation}\label{(15.2)}
    \left\{ \Ln z \right\} = \ln\left| z \right| + i \Arg z
\end{equation}
\begin{equation}\label{(15.3)}
    \Arg z= \left\{ \argt z + 2 \pi k \mid k \in \ZZ \right\}
\end{equation}
Из теоремы $2 \ \S 5$ (следствия $1$) можем видеть, что на $\CC \setminus \left(
    -\infty; 0 \right]$ существуют регулярные ветви:
\begin{equation}\label{(15.4)}
    \begin{split}
        & g_0(z) = \sqrt[n]{\left| z \right|}\exp\left( \frac{i}{n} \argm z \right) \\
        & h_0(z) = \ln\left| z \right| + i \argm z \\
        & \argm z \in \left( -\pi; \pi \right)
    \end{split}
\end{equation}
Это \textbf{главные регулярные ветви $\left\{\sqrt[n]{z}\right\}$ и $\Ln z$}.
\\
Рассмотрим \textbf{простейший случай}.
\\
Фиксируем $\varphi_0 \in \left( -\pi; \pi \right)$, $n \geq 2$.
\\
Пусть область
\begin{align*}
  & G_{1, \varphi_0} = \left\{ z \neq 0 \mid z = re^{i\varphi}, \ r > 0, \ \varphi \in \left( \frac{\varphi_0}{n}; \frac{\varphi_0 + 2\pi}{n} \right) \right\}
\end{align*}
$w = z^n$ однолистна на $G_{1, \varphi_0} \mapsto \CC \setminus
\lambda_{\varphi_0}$, гле $\lambda_{\varphi_0} = \left\{ z \mid z = re^{i
      \varphi_0}, \ r> 0\right\}\cup \{0\}$
\\
По следствию $1$ $\S 5$
\begin{equation}\label{(15.5)}
    \begin{split}
        & g_*(z) = \sqrt[n]{\left| z \right|}\exp\left( \frac{i}{n} \arg_* z \right) \\
        & \arg_* z \in \left( \varphi_0; \varphi_0+2\pi \right)
    \end{split}
\end{equation}
дает регулярную ветвь.
Пусть область
\begin{align*}
  & G_{2, \varphi_0} = \left\{ z \mid \Img z \in \left( \varphi_0; \varphi_0 + 2\pi \right) \right\}
\end{align*}
$w = e^z$ однолистна на $G_{2, \varphi_0}$, и
\begin{equation}\label{(15.6)}
    \begin{split}
        & h_*(z) = \ln\left| z \right| + i \arg_* z \\
        & \arg_* z \in \left( \varphi_0; \varphi_0+2\pi \right)
    \end{split}
\end{equation}
дает регулярную ветвь.
\lemma
Пусть $G$~--- область в $\CC \setminus \lambda_{\varphi_0}$. Тогда все
непрерывные ветви $\left\{ \sqrt[n]{z} \right\}$ и $\Ln z$ в $G$ являются
регулярными и удовлетворяют соотношениям:
\begin{equation}\label{(15.7)}
    g_k(z) = g_*(z) \exp\left( \frac{i}{n} 2 \pi k\right), \ k \in \{0, \dots, n-1\}
\end{equation}
\begin{equation}\label{(15.8)}
    h_k(z) = h_*(z) + 2i \pi k, \ k \in \ZZ
\end{equation}
\pr
Пусть $g(z)$~--- непрерывная ветвь $\left\{\sqrt[n]{z}\right\}$ в $G$; $g^n(z) \equiv
    z$, $g_*^n(z) \equiv z$. Деля,
    \begin{align*}
      & \left( \frac{g^n(z)}{g_*^n(z)} \right)^n = 1 \Rightarrow \frac{g(z)}{g_*(z)} \in \left\{ \sqrt[n]{1} \right\} = \exp \left( \frac{i}{n} 2 \pi k(z) \right), \ k(z) \in \left\{ 0, \dots, n-1 \right\}
    \end{align*}
В левой части функция непрерывная, в правой~--- ступенчатая, а значит, $k(z) =
k_* = const$, откуда видно, что $g(z)$ удовлетворяет \eqref{(15.7)}.
\\
Пусть $h(z)$~--- непрерывная ветвь $\Ln z$ в $G$; $e^{h(z)} \equiv z$,
$e^{h_*(z)} \equiv z$. Деля,
\begin{align*}
  & e^{h(z)-h_*(z)} = 1 \Rightarrow h(z) - h_*(z) = 2i\pi k(z)
\end{align*}
Значит, $k(z) = k_* = const$, откуда видно, что $h(z)$ удовлетворяет
\eqref{(15.8)}.
\lemma
В любом кольце $K_{r,R}(0) = \left\{ z \mid r < \left| z \right| < R \leq
    +\infty \right\}$ нет непрерывных ветвей $\left\{ \sqrt[n]{z} \right\}$ или
$\Ln z$.
\pr
Предположим противное: $\exists K_{r, R}(0)$: $\left\{ \sqrt[n]{z} \right\}$
имеет в нем непрерывную ветвь $g(z)$. 
\\
Тогда рассмотрим область $G = K \setminus \left( -\infty;0 \right]\subseteq \CC
\setminus \left( -\infty; 0 \right]$. Эта область удовлетворяет условиям леммы
$1$, но по этой лемме $\exists k_* \in \left\{ 0, \dots, n-1 \right\}$: $g(z) =
g_0(z)\exp \left( \dst \frac{i}{n} 2 \pi k_* \right)$, $z \in G$. Заметим, что
$\forall x \in \left( -R;-r \right)$
\begin{align*}
& g(x+i0) = \sqrt[n]{\left| z \right|}\exp \left( \frac{i \pi}{n} + \frac{2 i \pi k_*}{n}\right) \\
& g(x-i0) = \sqrt[n]{\left| z \right|}\exp \left( \frac{-i \pi}{n} + \frac{2 i \pi k_*}{n}\right) \\
\end{align*}
Они не равны, противоречие с предположением о непрерывности.
\\
Аналогично можно провести доказательство для $\Ln z$.
\corollary
Если $0 \in G$ или $\infty \in G$, то в $G$ не существует непрерывной ветви
$\left\{ \sqrt[n]{z} \right\}$ или $\Ln z$.
\\
Рассмотрим \textbf{общий случай}.
\\
Пусть $G \subseteq \CC$~--- односвязая область, $0 \not \in G$.
\lemma
В односвязной области $G$, $0 \not \in G$ существуют непрерывные ветви $\Arg z$.
\pr
Фиксируем $z_0 \in G$, $\psi_0 \in \Arg z_0$.
\\
Пусть $\gamma_{z_0z}$~--- гладкая кривая из $z_0$ в $z$, $ \gamma \subseteq G$.
Пусть
\begin{align*}
  & \varphi_{\gamma_{z_0z}} = \psi_0 + \Delta_{\gamma_{z_0z}}\argt z
\end{align*}
\begin{align*}
  & \varphi_{\gamma_{z_0z}} \in \Arg z
\end{align*}
\begin{align*}
  & \Delta_{\gamma_{z_0z}}\argt z = \Img \int_{\gamma_{z_0z}}\frac{d \zeta}{\zeta}
\end{align*}
Знаем, что $\dst \frac{1}{z}$ регулярна в $G$, а значит, по теореме Коши
\begin{align*}
  & \forall \os{\circ}{\gamma} \subseteq G \ \int_{\os{\circ}{\gamma}}\frac{d\zeta}{\zeta} = 0
\end{align*}
Тогда
\begin{align*}
  & \Phi(z) = \int_{\gamma_{z_0z}}\frac{d \zeta}{\zeta}
\end{align*}
регулярна в $G$;
\begin{equation}\label{(15.9)}
    \varphi_{\gamma_{z_0z}} = \varphi_0(z) = \psi_0 + \Delta_{\gamma_{z_0z}}\argt z \in \Arg z
\end{equation}
будет непрерывной, и, зная это, легко построить
\begin{equation}\label{(15.10)}
    \tilde{g}(z) = \sqrt[n]{\left| z \right|} \exp\left( \frac{i}{n} \varphi_0(z) \right)
\end{equation}
\begin{equation}\label{(15.11)}
    \tilde{h}(z) = \ln \left| z \right| + i \varphi_0(z)
\end{equation}
Это непрерывные ветви $\left\{ \sqrt[n]{z} \right\}$ и $\Ln z$ соответственно в
$G$.
\theorem
В односвязной $G$, $0 \not \in G$ существуют непрерывные ветви $\left\{
    \sqrt[n]{z} \right\}$ и $\Ln z$, являющиея регулярными ветвями и
удовлетворяюшщие соотношениям:
\begin{equation}\label{(15.12)}
    \tilde{g}_k(z) = \tilde{g}(z)\exp\left( \frac{2i\pi k}{n} \right), \ k \in \left\{ 0, \dots, n-1 \right\}
\end{equation}
\begin{equation}\label{(15.13)}
    \tilde{h}_k(z) = \tilde{h}(z) + 2 i \pi k, \ k \in \ZZ
\end{equation}
\pr
Всякая ветвь такова по лемме $1$. Докажем регулярность $\tilde{g}(z)$.
\\
Рассмотрим $z_1 \in G$: $\exists B_r(z_1) \subseteq G$, $0 \not \in B_r(z_1)$
$\Rightarrow$ $\exists \varphi_0 \in [-\pi, \pi): \ B_r(z_1) \subseteq \CC
\setminus \lambda_{\varphi_0}$.
\\
Значит, по лемме $1$ $\tilde{g}(z)$ регулярна в $B_r(z_1)$, а значит, и в $z_1$;
в силу произвольности выбора $z_1$ можем сказать, что функция регулярна в $G$.
\\
Аналогично можем провести доказательство для $\tilde{h}(z)$.

	\begin{flushright}
    \textit{Лекция 12 (от 13.10)}
\end{flushright}
\corollary
Пусть $G$ односвязна, $0 \not \in G$. Тогда для регулярной $h_k \in \Ln z$
справедлива формула Ньютона-Лейбница
\begin{equation}\label{(15.14)}
    h_k(z) = h_k(a) + \int_{\gamma_{az}}\frac{d\zeta}{\zeta}
\end{equation}
где $a \in G$, $\gamma_{az} \subseteq G$.
\pr
\begin{align*}
  & h_k(z) = h_0(z) + 2i \pi k
\end{align*}
\begin{align*}
  & h_0(z) = \ln \left| z \right| + i \left( \psi_0 + \Delta_{\gamma_{az}}\argt z \right)
\end{align*}
\begin{align*}
  & \Real \int_{\gamma_{az}} \frac{d\zeta}{\zeta} = \ln \left| z \right| - \ln \left| a \right|
\end{align*}
Действительно, пусть $z(t) = \gamma_{az}$, тогда
\begin{align*}
  & \Real \int_{\gamma_{az}} \frac{d\zeta}{\zeta} = \Real \int_0^1 \frac{xx'+yy'}{x^2+y^2} d \tau = \Real \int_0^1 \frac{d \sqrt{x^2+y^2}}{\sqrt{x^2+y^2}} = \ln \left| z(1) \right| - \ln \left| z(0) \right| = \ln \left| z \right| - \ln \left| a \right|
\end{align*}
\begin{align*}
  & h_k(z) = \ln \left| z \right| - \ln \left| a \right| + \ln\left| a \right| + i \left( \psi_0 + 2 \pi k \right) + i \Delta_{\gamma_{az}}\argt z = h_k(a) + \int_{\gamma_{az}}\frac{d\zeta}{\zeta}
\end{align*}
\section{$\S 16$. Регулярные ветви $\{\sqrt[n]{f(z)}\}$ и $\Ln f(z)$.}
\sug
$G$~--- область, $f$ регулярна на $G$, $\forall z \in G \ f(z) \neq 0$.
\begin{equation}\label{(16.1)}
    \left\{ \sqrt[n]{f(z)} \right\} = \sqrt[n]{\left| f(z) \right|} \exp \left( \frac{i}{n}\Arg f(z) \right)
\end{equation}
\begin{equation}\label{(16.2)}
    \Ln f(z) = \ln \left| f(z) \right| +i \Arg f(z)
\end{equation}
\begin{equation}\label{(16.3)}
    \Arg f(z) = \left\{ \argt f(z) + 2 \pi k \mid k \in \ZZ \right\}
\end{equation}
\Def
Пусть $\gamma$~--- непрерывная кривая в $G$, $f$ удовлетворяет предположению
$1$. Пусть $z(t)$~--- параметризация $\gamma$, $t \in [0;1]$. Пусть $\Gamma =
f(\gamma): \ w = f(z(t)), \ t \in [0;1]$. Тогда \textbf{приращением аргумента
  $f(z)$ вдоль кривой $\gamma$} называется
\begin{equation}\label{(16.4)}
    \Delta_{\gamma}\argt f(z) = \Delta_\Gamma\argt w = \Delta_{[0;1]}\argt f(z(t))
\end{equation}
\lemma
Пусть $f$, $f_1$, $f_2$ удовлетворяют предположению $1$ в области $G$. Тогда
\begin{enumerate}
    \item для любой непрерывной $\gamma \subseteq G$ выполняется логарифмическое
    свойство:
    \begin{equation}\label{(16.5)}
        \Delta_{\gamma}\argt (f_1f_2) = \Delta_\gamma\argt f_1 + \Delta_{\gamma}\argt f_2
    \end{equation}
    \item если $\gamma \subseteq G$ разбита точкой $A \in \gamma$ на части
    $\gamma_1$, $\gamma_2$, то
    \begin{equation}\label{(16.6)}
        \Delta_{\gamma}\argt f = \Delta_{\gamma_1}\argt f + \Delta_{\gamma_2}\argt f
    \end{equation}
    \item если $\gamma$~--- кусочно гладкая кривая в $G$, то
    \begin{equation}\label{(16.7)}
        \Delta_{\gamma}\argt f(z) = \Img \int_\Gamma \frac{dw}{w} = \Img \int_{\gamma} \frac{f'(\zeta)}{f(\zeta)}d\zeta
    \end{equation}
\end{enumerate}
\pr
Утверждения леммы очевидны из определения приращения аргумента функции.
\lemma
Пусть $G$ односвязна, $\os{\circ}{\gamma} \subseteq G$ замкнута и непрерывна.
Пусть $f$ удовлетворяет предположению $1$. Тогда
\begin{equation}\label{(16.8)}
    \Delta_{\os{\circ}{\gamma}}\argt f(z) = 0
\end{equation}
\pr
Пусть $\os {\circ}{\gamma} \subseteq G$ гладкая и замкнутая, параметризованная
$z(t)$.
\\
$\forall z_0 \in \os{\circ}{\gamma}$ существует непрерывная деформация $z(t,
\alpha)$ такая, что $z(t,a) = z(t)$, $z(t,b) = z_0$. По теореме $3$ $\S 14$
$I(\alpha) = const$, а значит,
\begin{align*}
  \Delta_{[0;1]}\argt f(z) = \Delta_{[0;1]}\argt f(z_0) = 0
\end{align*}
\lemma
Пусть $f$ удовлетворяет предположению $1$ в области $G$. Если в области $G$
сществуют регулярные ветви $h_0(z)$ или $g_0(z)$~--- регулярные ветви $\Ln f(z)$
или $\left\{ \sqrt[n]{f(z)} \right\}$, то все их непрерывные ветви регулярны и
удовлетворяют соотношениям:
\begin{equation}\label{(16.9)}
    h_k(z) = h_0(z) + 2i \pi k, \ k \in \ZZ
\end{equation}
\begin{equation}\label{(16.10)}
    g_k(z) = g_0(z) \exp\left( \frac{i}{n} 2 \pi k\right), \ k \in \{0, \dots, n-1\}
\end{equation}
\pr
Доказательство полностью аналогично доказательству леммы $15.1$.
\lemma
Пусть $f$ удовлетворяет предположению $1$ в $G$. Если в $G$ существует
регулярная ветвь $h(z)$ многозначной функции $\Ln z$, то $\forall a,b \in G$
\begin{equation}\label{(16.11)}
    h(b) = h(a) + \ln \left| \frac{f(b)}{f(a)} \right| + i \Delta_{\gamma_{ab}}\argt f(z)
\end{equation}
\pr
\begin{align*}
  & h(z) = \ln \left| f(z) \right| + i \Img h(z)
\end{align*}
В силу регулярности $h$ $\Img h$ будет гармонической функцией от $x, y$. Значит,
\begin{align*}
  & \varphi(z) = \Img h(z) \in \Arg f(z)
\end{align*}
Пусть $a \in G$, $\psi_0 \in \Arg f(a)$, $h(a) = \ln \left| f(a) \right| + i
\psi_0$. Пусть $\gamma_{az}$~--- гладкая кривая с параметризацией $z(t)$, $t \in
[0;1]$, $\varphi(z(t))$~--- гладкая ветвь $\Arg f(z(t))$.
\begin{align*}
  & \Delta_{\gamma_{az}}\argt f(z) = \varphi(z(t))-\varphi(z(0)) = \Img h(z) - \Img h(a)
\end{align*}
\begin{align*}
  & h(b) = \ln \left| f(b) \right| + i \Img h(b)
\end{align*}
\begin{align*}
  & h(a) = \ln \left| f(a) \right| + i \Img h(a)
\end{align*}
\begin{align*}
  & h(b) - h(a) = \ln \left| \frac{f(b)}{f(a)} \right| + i \left( \Img h(b) - \Img h(a) \right)
\end{align*}
Отсюда следует \eqref{(16.11)}.
\theorem
Пусть $f$ удовлетворяет предположению $1$ в $G$. Тогда у  $\Ln f(z)$ существуют
регулярные ветви в $G$ тогда и только тогда, когда выполняется \eqref{(16.8)}
для любой замкнутой $\os{\circ}{\gamma}\subseteq G$.
\pr
Докажем в обе стороны.
\begin{itemize}
    \item Необходимость.
    \\
    Если есть регулярные ветви $h(z) \in \Ln z$  $G$, то из \eqref{(16.11)} при
    $a = b$ получаем \eqref{(16.8)}.
    \item Достаточность.
    \\
    Пусть выполнено \eqref{(16.8)}. Фиксируем $a \in G$, $h(a) \in \Ln f(a)$.
    Рассмотрим
    \begin{equation}\label{(16.12)}
        h(z) = h(a) + \ln \left| \frac{f(z)}{f(a)} \right| + i \Delta_{\gamma_{az}}\argt f(z)
    \end{equation}
    В действительности приращение аргумента может зависеть от $\gamma$, поэтому
    корректнее расссматривать это как $h(\gamma_{az})$.
    \\
    По \eqref{(16.8)}
    \begin{align*}
      & \Delta_{\gamma_{az}}\argt f(z) = \Delta_{\tilde{\gamma}_{az}}\argt f(z)
    \end{align*}
    \begin{align*}
      & \os{\circ}{\gamma} = \gamma_{az}+\tilde{\gamma}^{-1}_{az}
    \end{align*}
    Значит, $h(z)$~--- функция лишь точки $z$.
    \begin{equation}\label{(16.13)}
        \exists \psi_0 \in \Arg f(a): \ h(a) = \ln \left| f(a) \right| + i \psi_0
    \end{equation}
    Из \eqref{(16.12)} следует, что
    \begin{align*}
      & h(z) = \ln \left| f(z) \right| + i \left( \psi_0 + \Delta_{\gamma_{az}}\argt f(z) \right) \in \Ln f(z)
    \end{align*}
    \begin{align*}
      & \Real \int_{\gamma_{az}}\frac{f'(\zeta)}{f(\zeta)}d\zeta = \ln \left| f(z) \right| - \ln \left| f(a) \right|
    \end{align*}
    \begin{align*}
      & \Delta_{\gamma_{az}}\argt f(z) = \Img \int_{\gamma_{az}}\frac{f'(\zeta)}{f(\zeta)}d\zeta
    \end{align*}
    Отсюда следует, что
    \begin{align*}
      & h(z) = \ln \left| f(a) \right| + \ln \left| f(z) \right| - \ln \left| f(a) \right| + i \Delta_{\gamma_{az}} \argt f(z) + i \psi_0 = \ln \left| f(a) \right| + \int_{\gamma_{az}}\frac{f'(\zeta)}{f(\zeta)}d \zeta +  i \psi_0 
    \end{align*}
    Функция
    \begin{align*}
      & \varphi(z) = \int_{\gamma_{az}}\frac{f'(\zeta)}{f(\zeta)}d\zeta
    \end{align*}
    регулярна по теореме $3$ $\S 6$, а значит, и $h(z)$ регулярна.
\end{itemize}
\lemma
Пусть $f$ удовлетворяет предположению $1$ в $G$, $n \in \NN$, $n \geq 2$. Если в
$G$ существует регулярная ветвь $g(z)$ многозначной функции $\left\{
    \sqrt[n]{f(z)} \right\}$, то $\forall a, b \in G$
\begin{equation}\label{(16.14)}
    g(b) = g(a) \sqrt[n]{\left| \frac{f(b)}{f(a)} \right|}\exp \left( \frac{1}{n}\Delta_{\gamma_{ab}}\argt f(z) \right)
\end{equation}
\pr Доказательство разбиваем на два этапа.
\begin{enumerate}
    \item Пусть для начала $G$ односвязна.
    \\
    По лемме $2$ $\forall \os{\circ}{\gamma} \subseteq G$ выполняется
    \eqref{(16.8)} и тогда, по лемме $4$:
    \begin{align*}
      & \exists h(z) = h(a) + \ln \left|\frac{f(z)}{f(a)}\right|+ i\left(\Delta_{\gamma_{az}}\argt f(z)\right) \in \Ln f(z)
    \end{align*}
    регулярные ветви.
    \begin{align*}
      & g_0(z) = \exp \left(\frac{i}{n} h(z)\right)
    \end{align*}
    \begin{align*}
      & \exists \psi_0 \in \Arg f(a): \ h(z) = \ln \left| f(z) \right| + i \left( \psi_0 + \Delta_{\gamma_{az}}\argt f(z) \right)
    \end{align*}
    Тогда
    \begin{align*}
      & g_0(z) = \sqrt[n]{\left| f(z) \right|}\exp \left( \frac{i}{n} \left( \psi_0 + \Delta_{\gamma_{az}}\argt f(z) \right) \right)\in \left\{ \sqrt[n]{f(z)} \right\}
    \end{align*}
    регулярная ветвь.
    \\
    По лемме $3$ $\exists k \in \left\{ 0, \dots, n-1 \right\}: \ g(z) = g_0(z)$.
    Отсюда следует \eqref{(16.14)}.
    \item $G$~--- область.
    \begin{align*}
      &\forall a, b \in G, \ \forall \gamma_{ab} \subseteq G
    \end{align*}
    можем выбрать такие $a=z_0, z_1, \dots, z_{K-1}, z_K = b \in \gamma_{ab}$,
    что $\left| z_k - z_{k-1} \right|< l_k$, где $l_k$~--- расстояние до
    границы.
    \begin{align*}
      & \exists \left\{ B_\varepsilon(z_k) \right\}_{k=0}^K \subseteq G: \ z_{k+1} \in B_{\varepsilon}(z_k)
    \end{align*}
    Каждый из $B_{\varepsilon}(z_k)$ односвязен, поэтому выполняется
    \eqref{(16.14)}:
    \begin{align*}
      & \forall k \in \left\{ 0, \dots, K-1 \right\} \ \frac{g(z_{k+1})}{g(z_k)} = \sqrt[n]{\frac{f(z_{k+1})}{f(z_k)}} \exp \left( \frac{i}{n} \Delta_{\gamma_{z_kz_{k+1}}}\argt f(z) \right)
    \end{align*}
    Перемножая по всем $k$, получаем
    \begin{align*}
      & \frac{g(b)}{g(a)} = \sqrt[n]{\left| \frac{f(b)}{f(a)} \right|}\exp \left( \frac{i}{n} \sum_{k=0}^{K-1} \Delta_{\gamma_{z_kz_{k+1}}} \argt f(z) \right) = \sqrt[n]{\frac{\left| f(b) \right|}{\left| f(a) \right|}}\exp \left( \frac{i}{n}\Delta_{\gamma_{ab}}\argt f(z) \right)
    \end{align*}
    Это ни что иное, как \eqref{(16.14)}.
\end{enumerate}\theorem
Пусть $f$ удовлетворяет условиям предположения $1$ в $G$. Тогда существуют
регуляртые ветви $\left\{ \sqrt[n]{f(z)} \right\}$тогда и только тогда, когда
\begin{equation}\label{(16.15)}
    \forall \os{\circ}{\gamma} \subseteq G \ \exists k(\os{\circ}{\gamma}) \in \ZZ: \ \Delta_{\os{\circ}{\gamma}}\argt f(z) = 2\pi n k(\os{\circ}{\gamma}))
\end{equation}
\end{document}